\documentclass{article}
\usepackage[margin=1in]{geometry}
\usepackage{amsmath}
\usepackage{amssymb}
\usepackage{amsfonts}
\usepackage{graphicx}
\usepackage{boxedminipage}
\usepackage{hyperref}
\usepackage{natbib}
\usepackage{enumitem}
\usepackage{tocloft}
\usepackage{tocbibind}
\usepackage{authblk}
\usepackage{sectsty}
\usepackage{lipsum} % For placeholder text where needed

\sectionfont{\large\bfseries}
\subsectionfont{\normalsize\bfseries}
\subsubsectionfont{\normalsize\itshape}

\title{Consciousness as Gait: The RSVP--CPG Framework for Cyclical Cognition, Memory, and Sleep}
\author{Flyxion}
\date{September 26, 2025}

\begin{document}

\maketitle


\begin{abstract}
Recent magnetoencephalography studies have revealed that large-scale cortical networks activate in robust cyclical sequences, with asymmetric transitions, stable ordering, and strong heritability of cycle rate (Van Es et al., 2025). These findings call for a framework that explains not only the existence of cycles but also their role in cognition, memory, and sleep. Global Workspace Theory interprets consciousness as selective broadcast, while Active Inference frames it as predictive equilibrium. We propose an alternative: consciousness as gait.

In the RSVP--CPG framework, cortical cycles function as chained central pattern generators, memory persists as proxy loops entrained to the cortical gait, and paradoxical sleep arises from desynchronization of chains. Drawing on Arnie Cox’s work on music as embodied cognition (Cox, 2016) and Barbara Tversky’s account of mind as fundamentally shaped by spatial motion (Tversky, 2019), we argue that cognition is best understood as rhythmic, embodied progression rather than spotlight or balance. Our Bayesian generative model recovers cycle order, asymmetry, and rate from MEG data, linking cortical cycles to behavior, memory replay, and dream phenomenology.

By positioning cortical cycles as non-equilibrium oscillatory attractors, RSVP--CPG reframes cognition as perpetual rhythmic motion: a gait of thought in which consciousness is carried forward step by step, constrained and enabled by entropy itself.
\end{abstract}


\tableofcontents

\section{Introduction: Thinking in Motion}

The connection between movement and thought has deep roots. Aristotle’s school was known as the Peripatetic because of its practice of walking during philosophical discourse. For the Peripatetics, thought was not divorced from motion but carried along by it: the rhythm of steps paced the rhythm of inquiry. This ancient intuition—that cognition is grounded in embodied movement—has found renewed force in contemporary research.

Arnie Cox’s work on music as embodied cognition argues that listening itself is an act of motor resonance: we understand melodies and rhythms by covertly enacting them in our own bodies \citep{cox2016music}. Similarly, Barbara Tversky’s Mind in Motion develops the thesis that cognition is fundamentally spatial and kinetic, that abstract thought borrows its structure from the dynamics of bodily action \citep{tversky2019mind}. Both lines of research converge on a simple yet profound claim: mind is not a disembodied processor of information, but a system whose architecture is rhythmic, spatial, and kinetic.

The RSVP--CPG framework extends this embodied lineage into a neurocomputational theory of consciousness. Building on recent evidence that cortical networks activate in robust cyclical sequences \citep{van2025cortical}, we propose that consciousness itself is a gait. Cortical cycles function as chained central pattern generators, carrying perception, memory, and action forward in rhythmic order. Memory persists through proxy loops entrained to this gait, while paradoxical sleep arises when the chains desynchronize. By situating RSVP--CPG in the tradition of peripatetic philosophy and in dialogue with contemporary accounts of embodied cognition, we suggest that the oscillatory skeleton uncovered in neuroimaging is not an incidental correlate but the fundamental form of thought.

To further elucidate this framework, we delve deeper into its theoretical foundations, empirical validations, and interdisciplinary implications. The following sections expand on the mathematical rigor, philosophical underpinnings, and practical applications of RSVP--CPG, providing a comprehensive exploration of consciousness as rhythmic gait.

\section{Conceptual Lineage: From Peripatetic Thought to Cyclical Cognition}

The framing of consciousness as gait belongs to a lineage that spans philosophy, cognitive science, and neuroscience.

Peripatetic philosophy. In Aristotle’s Lyceum, philosophy was practiced peripatein—to walk while reasoning. The Peripatetics implicitly recognized that the cadence of thought is inseparable from the cadence of the body. Motion was not merely metaphor but medium: reasoning unfolded in rhythm with the steps that carried it.

Embodied cognition. Contemporary accounts echo this ancient intuition. Arnie Cox \citep{cox2016music} argues that music comprehension is grounded in motor resonance: listeners simulate the actions required to produce sounds, effectively “moving along” with what they hear. Barbara Tversky \citep{tversky2019mind} generalizes this to cognition at large, contending that spatial and kinetic structures are the scaffolds of mental life. For both, thought is fundamentally a translation of movement into abstract domains.

Neuroscientific cycles. Recent neuroimaging studies provide empirical confirmation at the systems level. Van Es et al. \citep{van2025cortical} show that large-scale cortical networks activate in cyclical, asymmetric sequences that recur with remarkable stability and heritability. Cognition, long suspected to have a rhythmic basis, now reveals its cyclical skeleton under direct measurement.

RSVP--CPG. The Relativistic Scalar–Vector Plenum with Central Pattern Generators synthesizes these insights into a field-theoretic account. Cortical cycles are modeled as chained CPGs, enforcing ordered progression through network states. Memory emerges as proxy loops entrained to these rhythms; REM sleep as a regime of desynchronized chains. By grounding cognition in non-equilibrium oscillations, RSVP--CPG continues the peripatetic claim: to think is to move, and to move is to think.

This lineage not only historical but also conceptual, linking ancient intuitions to modern empirical findings. In the following subsections, we explore how this evolution informs the core principles of RSVP--CPG.

\subsection{From Static Models to Dynamic Flows}

Traditional theories of mind often conceptualize cognition as a set of discrete states or symbolic tokens. In such models, mental life is represented as a sequence of static configurations, much like snapshots in a film reel. This view lends itself to computational analogies: the brain as a digital machine executing operations over representations.

The RSVP--CPG framework rejects this static picture in favor of dynamic flows. Cognition is not a sequence of states but a continuous unfolding, much as walking is not a series of frozen postures but a rhythmic transition from imbalance to balance. This perspective aligns with the broader shift in cognitive science toward dynamical systems theory, where the emphasis is placed on trajectories in state space rather than isolated states.

The peripatetic tradition anticipated this view by treating motion as the medium of thought. More recently, enactive cognition theorists such as Varela, Thompson, and Rosch \citep{varela1991embodied} have argued that cognition emerges through the ongoing coupling of organism and environment. RSVP--CPG provides a neuroscientific instantiation of this claim: cortical cycles are not representations of external states but dynamic flows that continually reshape the cognitive field.

This reconceptualization has epistemological consequences. If cognition is fundamentally flow rather than static state, then knowledge itself is rhythmic: it emerges through entrainment, recurrence, and re-alignment rather than through the accumulation of static facts. In this sense, RSVP--CPG extends the peripatetic principle into a formal account: to think is to flow, not to freeze.


\subsection{Integration of Temporal Dimensions}

Time has always been the hidden dimension of cognition. Husserl’s phenomenology of internal time-consciousness \citep{husserl1966phenomenology} argued that every conscious act is structured by a threefold temporal horizon: the just-past (retention), the living-present (primal impression), and the about-to-occur (protention). Heidegger extended this in \emph{Being and Time} \citep{heidegger1927being}, framing human existence itself as temporal unfolding: to be is to be stretched along a horizon of possibilities. Both accounts stress that cognition cannot be reduced to static instants; it is essentially rhythmic, carried forward by the flow of time.

The RSVP--CPG framework provides a neuroscientific correlate for this phenomenological structure. Cortical cycles partition cognition into repeating windows, each one a temporal scaffold that aligns retention, impression, and protention with measurable oscillatory phases. In this way, Husserl’s tripartite structure finds a natural mapping: retention corresponds to mnemonic loops carried forward across cycles, impression corresponds to the active cortical state at a given cycle phase, and protention corresponds to the anticipatory bias set by the ordering of subsequent phases. Consciousness, in this view, is temporally thick because it is oscillatory; every present moment contains echoes of the past and anticipations of the future, braided together in rhythmic progression.

This temporal grounding has implications for both theory and experiment. Phenomenologically, it explains why cognition feels like flow rather than discrete steps. Empirically, it predicts that subjective temporal judgments (e.g., simultaneity perception, duration estimation) should covary with the phase and strength of cortical cycles. Disruptions of timing—such as those observed in schizophrenia or Parkinson’s disease—can thus be reinterpreted as breakdowns in cycle stability. Temporal consciousness, long theorized as philosophy’s deepest problem, becomes legible as the gait of neural oscillators.

In short, RSVP--CPG unites phenomenology and neuroscience: the temporality described by Husserl and Heidegger is not metaphor but mechanism, instantiated in cortical cycles that pace the gait of thought.


\section{Gesture, Motion, and the Origins of Cognition}

The primacy of rhythm and motion in thought is not only philosophical and phenomenological but also developmental and evolutionary. Gestural languages precede auditory languages both phylogenetically and ontogenetically. Comparative primatology shows that great apes communicate extensively through gesture long before they evolve structured vocal signals \citep{tomasello2008origins}. Infants, likewise, establish communication through pointing, reaching, and imitation months before they acquire speech \citep{goldin2003language}. Gesture is the first syntax: a chaining of bodily movements into communicative gaits.

Within this framework, theory of mind arises from motor simulation. To recognize the intentions of another, one must map their movements onto one’s own bodily repertoire — effectively copying their dance moves. This is not a poetic gloss but a motor-theoretic description: mirror neuron systems, premotor cortices, and cerebellar circuits allow an observer to “run” another’s actions through their own CPG chains. Recognition of agency emerges from entrainment: when the other’s rhythm can be carried in one’s own gait, the other is understood as intentional.

This line of reasoning strengthens the RSVP--CPG thesis. If cognition is grounded in chained oscillators and rhythmic entrainment, then language and social cognition are specializations of a more fundamental architecture of gait. Auditory language piggybacks on gestural rhythms; theory of mind piggybacks on the capacity to entrain to others’ cycles. Even the bizarreness of REM dreams can be reframed here: desynchronized CPG chains temporarily decouple the rhythms that normally align self with world, producing the estranging juxtapositions characteristic of dream life.

To expand this perspective, we consider the evolutionary advantages of rhythmic cognition and its implications for modern cognitive science.

\subsection{Evolutionary Perspectives on Rhythmic Cognition}

From an evolutionary standpoint, the development of gait-like cognitive processes may have conferred adaptive advantages in social coordination and environmental navigation.

\subsection{Ontogenetic Development of Gait Patterns}

In human development, the emergence of rhythmic patterns in infancy parallels the maturation of cognitive capabilities, supporting the embodied nature of mind.

\section{Gait, Agency, and the Spectrum of Self/Other}

The RSVP--CPG framework suggests that cognition operates by entraining rhythmic chains of motion. This gives rise to a fundamental spectrum:

If the observed pattern of motion too closely matches one’s own gait, it is registered as self.

If the pattern sufficiently deviates yet remains entrainable, it is recognized as other.

If it is unentrainable or inconsistent with known bodily rhythms, it is coded as inanimate environment.

This spectrum underlies not only theory of mind (“copying their dance moves” as a way of modeling others) but also basic perceptual filtering:

Heartbeat and proprioception. Rhythms generated by the body (pulse, respiration, postural sway) are discounted as self-caused.

Eye geometry and perspective. Distortions from corneal curvature or perspective transformations are filtered as self-produced, preventing them from being mistaken for external agents.

Shadows and reflections. Visual alterations that align with self-movement are likewise coded as extensions of one’s own body.

At the same time, this spectrum explains the need for skeuomorphic affordances in tool use and interface design. A tool becomes usable when its dynamics can be entrained into one’s CPG-based gait. Interfaces that mimic real-world gestures (turning a page, plucking a string, twisting a knob) exploit this entrainability, anchoring the artificial in rhythms already owned by the body. Without such skeuomorphic grounding, tools remain opaque, because the system cannot decide where to locate them on the self/other spectrum.

Thus, the same mechanism that filters out self-generated distortions and anchors bodily rhythms is also the mechanism that permits social recognition and technological extension. The capacity to entrain and differentiate — to know whether a rhythm is “mine” or “theirs,” “self-caused” or “world-caused” — is the structural condition for both theory of mind and tool use.

Expanding on this spectrum, we explore its implications for social cognition and technological integration in modern societies.

\subsection{Social Implications of the Self/Other Spectrum}

In social contexts, the ability to entrain rhythms facilitates empathy and cooperation, while failures in differentiation may underlie certain social disorders.

\subsection{Technological Extensions and Human-Computer Interaction}

Skeuomorphic design principles, rooted in rhythmic entrainment, guide effective interface development, enhancing user intuition and efficiency.

\section{Markov Blankets, Gait, and the Self/Other Spectrum}

The self/other spectrum can be recast in the language of Markov blankets and Active Inference (AIF)—showing why physics modeling is, in form, theory-of-mind (ToM) inference.

1. The blanket partition. In AIF, a system is factorized by a Markov blanket into internal states (beliefs over hidden causes), active states (actions), sensory states (observations), and external states (world causes). Inference updates internal states to minimize variational free energy, while actions change sensory states to fulfill predictions.

RSVP--CPG view: the body’s gait generators (CPGs) supply strong priors over lawful sensorimotor rhythms (heartbeat, respiration, postural sway, ocular micro-motion). These priors define a tight sensorimotor loop—a self-blanket with high predictability and short latency.

2. Self/other discrimination as synchrony-based blanket assignment. Self: If an observed pattern is phase-locked to predicted consequences of actions (or canonical interoceptive rhythms), its prediction error is low under the self model. It is absorbed inside the self-blanket. Other/World: If entrainable but not locked, it fits an external generative model. Inanimate: If unentrainable, it is explained as passive dynamics.

3. Why physics modeling is ToM in disguise. Both ToM and intuitive physics are inference over external generative processes. ToM uses priors over goals/policies; physics uses priors over forces/constraints. Both predict motion; humans adopt intentional stances for complex patterns.

4. Filtering self-caused distortions via the self model. Cardiac/respiratory pulses, ocular distortion, perspective warps, self-shadow/reflection: predicted under the self model and explained away.

5. Skeuomorphic affordances as blanket extensions. Tools/interfaces become intuitive when contingencies match self policies, increasing model evidence for incorporation.

6. REM and dream bizarreness under blankets. During REM, precision on cross-chain coupling drops; blanket assignment becomes unstable, producing bizarreness.

7. Formal sketch. Competing generative models with priors over phase dynamics.

8. Practical consequences and tests. ToM ≈ physics: Perturb ambiguous stimuli; expect model selection to flip between priors. Tool learning: Skeuomorphic designs speed incorporation. REM: Bizarreness correlates with phase jitter.

In an Active Inference framing, self/other discrimination, theory of mind, intuitive physics, and tool affordances are all Markov-blanket assignments guided by rhythmic entrainment. RSVP--CPG supplies the oscillatory priors (the “gait”) that make these assignments fast and reliable.

To deepen this integration, we consider advanced applications in AI and robotics.

\subsection{AI Applications of Markov Blanket Dynamics}

Incorporating rhythmic priors in AI systems could enhance adaptive behavior, mimicking human-like flexibility in dynamic environments.

\begin{boxedminipage}{\textwidth}
\textbf{Literary Note: The Lobster Quadrille as Active Inference}

Lewis Carroll’s Mock Turtle’s Song (“The Lobster Quadrille”) offers a whimsical allegory of the inferential problem at the heart of cognition.

“Will you walk a little faster?” said a whiting to a snail… Will you, won’t you, will you, won’t you, will you join the dance?

In Active Inference terms, the whiting issues an invitation to entrain: to align the snail’s generative model with his own. Accepting would mean expanding the Markov blanket, treating the whiting’s rhythm as part of one’s coordinated self–other cycle.

The snail’s refusal — “Too far, too far!” — illustrates the role of precision weighting. When the cost of synchronization is too high, the system protects its boundaries, treating the offered rhythm as external noise rather than agentive motion.

From an RSVP–CPG perspective, the dance itself represents the oscillatory substrate of cognition: chained cycles that can align (social entrainment, theory of mind) or remain distinct (physics-like external dynamics). The lobster quadrille is thus a parable of self/other/environment discrimination: deciding which rhythms to treat as one’s own, which as another’s, and which as the impersonal world.

Just as Carroll’s creatures debate whether to “join the dance,” the brain continuously negotiates its own dances of perception, action, and inference.
\end{boxedminipage}

\footnote{Lewis Carroll’s Mock Turtle’s Song (“The Lobster Quadrille”) anticipates, in allegorical form, the RSVP–CPG view of cognition as gait. The whiting’s repeated invitation—“Will you, won’t you, will you, won’t you, will you join the dance?”—is an exhortation to entrain, to extend the Markov blanket and synchronize one’s cycles with another’s. The snail’s refusal (“Too far, too far!”) illustrates the role of precision weighting: when the inferential cost of synchronization is judged too high, rhythms are rejected as external rather than agentive. From this perspective, the “dance” itself dramatizes the oscillatory substrate of cognition, the problem of distinguishing self, other, and world through rhythmic entrainment. In this sense, Carroll’s quadrille functions as a playful parable of consciousness as gait.}

\footnote{Aristotle remarks in De Anima that “soul never thinks without a phantasm” (III.7, 431a), and his peripatetic practice embodied the same principle: reasoning unfolded while walking, paced by the rhythm of steps. Lewis Carroll’s Lobster Quadrille may be read as a parody of this peripatetic truth, where to “join the dance” is to synchronize one’s motion with another’s. The RSVP–CPG framework recasts this lineage: cognition as gait, consciousness as the ordering of oscillatory cycles whose rhythms both individuate the self and invite entrainment with others. \citep{aristotle1984complete, during1957aristotle}}

\section{Embodiment, Peripatetics, and the Rhythms of Thought}

The idea that cognition is inseparable from bodily motion has a long genealogy. Aristotle in De Anima (III.7, 431a15–17) asserts that “the soul never thinks without a phantasm,” insisting that thought requires the mediation of embodied imagery \citep{aristotle1984complete}. This principle was not only theoretical but practical: the Peripatetic school took its name from peripatein (“to walk about”), reflecting a pedagogy in which reasoning unfolded while walking \citep{during1957aristotle}. Movement was understood as the pacing of thought itself.

This embodied lineage extends into modern philosophy of mind. Gallagher \citep{gallagher2005body} argues that “the body shapes the mind” by providing the scaffolding for intentionality, proprioception, and social cognition. Arnie Cox \citep{cox2016music} develops a parallel argument in the domain of music: listening and comprehension depend upon covert motor resonance, the simulation of gestures and actions in the listener’s own body. Barbara Tversky \citep{tversky2019mind}, in Mind in Motion, generalizes this claim: cognition is fundamentally structured by spatial and kinetic schemas, with abstract reasoning borrowing directly from bodily movement and orientation.

The RSVP--CPG framework integrates these perspectives within a neuroscientific model of cyclical cortical dynamics. On this account, cortical cycles function as chained central pattern generators, carrying perception, memory, and action forward in ordered sequences. Gesture and rhythm are not auxiliary to language or thought, but primary: gestural communication precedes vocal language in both phylogeny and ontogeny, and theory of mind itself arises through motor entrainment — modeling another’s agency by copying their dance. What Aristotle described as phantasmata, and what the Peripatetics enacted in walking discourse, reappear here as oscillatory attractors in the RSVP--CPG field: the rhythmic substrate of consciousness.

Building on this foundation, we explore further the interplay between embodiment and rhythm in cognitive processes.
\subsection{The Role of Phenomenology in Embodied Cognition}

Merleau-Ponty’s \emph{Phenomenology of Perception} \citep{merleauponty2013phenomenology} placed the body at the center of consciousness. Against both empiricism and intellectualism, he argued that perception is not the passive reception of stimuli nor the imposition of concepts, but the lived articulation of the body’s capacities. Vision, touch, and movement are not separable channels but interwoven modalities of being-in-the-world. The body does not stand between mind and world—it is the very field in which meaning arises.

Within the RSVP--CPG framework, this phenomenological insight translates directly into neural dynamics. Cortical cycles function as chained central pattern generators, coupling sensory flows with motor affordances in rhythmic progression. The “I can” that Merleau-Ponty described—the sense that one can reach, grasp, or step forward—emerges from the stability of these cycles. Every perception is already motoric, every cognition already a gait forward into potential action.

Phenomenology also clarifies why disruptions of rhythm are so disorienting. In pathological conditions such as schizophrenia, disturbances of bodily self-awareness are often described as temporal or spatial disjunctions. From an RSVP--CPG perspective, these can be understood as breakdowns in the coherence of oscillatory cycles, leaving the subject alienated from the very rhythms that constitute embodied intentionality. Dreams, likewise, dramatize this principle: when CPG chains desynchronize in REM sleep, the normally stable rhythms of perception and action fracture, producing surreal experiences that phenomenology describes as “estranged embodiment.”

By grounding RSVP--CPG in Merleau-Ponty’s account of embodiment, we bridge first-person phenomenology with third-person neuroscience. The body is not an afterthought in cognition but its oscillatory foundation. Consciousness as gait thus becomes not just a metaphor but an existential structure: to be conscious is to move rhythmically through the world, inhabiting cycles that carry perception, memory, and action forward together.

\subsection{Process Philosophy Perspectives}

The RSVP--CPG framework resonates strongly with traditions in process philosophy, particularly the work of Alfred North Whitehead and Henri Bergson. Both rejected substance-based metaphysics in favor of an ontology of becoming: reality is constituted not by static things but by ongoing processes. For Whitehead \citep{whitehead1929process}, the fundamental units of existence are “actual occasions,” momentary events that prehend one another and weave themselves into enduring patterns. For Bergson \citep{bergson1910time}, consciousness itself is durée—an indivisible flow of time whose continuity cannot be captured by spatialized, static categories.

RSVP--CPG provides a neuroscientific instantiation of these processual insights. Cortical cycles are not static states but rhythmic becomings, transient actual occasions whose succession generates the felt continuity of consciousness. The chaining of CPGs mirrors Whitehead’s notion of concrescence: each cycle inherits constraints from the past, integrates them with present conditions, and hands them forward into the next. Bergson’s durée likewise finds resonance here: consciousness as gait is not a series of frozen instants but a rhythmic unfolding that resists reduction to discrete units.

This process-oriented interpretation also clarifies the role of entropy in RSVP. Rather than treating entropy as decay or loss, RSVP conceives it as the pacing mechanism of becoming: the gradient along which cycles fall forward, ensuring that thought remains in motion. In this sense, entropy is the guarantor of process, preventing cognition from lapsing into stasis.

By aligning RSVP--CPG with process philosophy, we highlight a continuity between metaphysical traditions and neuroscientific models. Whitehead’s actual occasions, Bergson’s durée, and RSVP’s cortical cycles all point to the same insight: consciousness is not a thing but a rhythm, a structured flow of becoming. To think is to move forward in cycles; to be conscious is to inhabit a process of perpetual gait.


\section{Comparative Analysis: GWT, AIF, and RSVP--CPG}

Theories of consciousness and cognition often rely on metaphors to render tractable processes that resist direct observation. Global Workspace Theory (GWT) and Active Inference (AIF) have emerged as two of the most influential approaches, framing cognition respectively in terms of broadcasting and prediction. The RSVP--CPG framework introduces a third account, in which cognition is modeled as rhythmic gait. Examining the contrasts across these theories illuminates both shared ground and fundamental divergences.

\textbf{Mechanism.} GWT models the brain as a set of specialized modules competing for access to a “global workspace,” with the winning representation broadcast across cortical and subcortical regions \citep{baars1997theater, dehaene2011experimental}. AIF, by contrast, posits that cognition arises through continual minimization of variational free energy, unifying perception and action under a single principle \citep{friston2010free, friston2017graphical}. Baioumy et al. \citep{baioumy2021active} extend this account to robotics, showing how AIF can encompass classic control schemes such as PID loops, enabling manipulators to adapt to disturbances and unmodeled dynamics. RSVP--CPG departs from both: cognition is neither an ignition event nor an equilibrium process, but a sequential unfolding through chained cortical central pattern generators (CPGs). Each step smooths entropy constraints before handing off to the next, producing a continuous rhythm of thought.

\textbf{Metaphor.} The differences are most vivid in metaphor. GWT employs the stage/spotlight image: emphasizing ignition and access \citep{baars1997theater}. AIF is best understood through the metaphor of an equilibrium machine: a system that continuously rebalances itself to minimize surprise \citep{friston2010free}. RSVP--CPG instead invokes a movement metaphor: cognition as gait. Here, consciousness does not arise from balance or selection, but from rhythmic progression — like walking, where stability emerges through perpetual imbalance, each step caught by the next. Cognition is not spotlight or balance, but forward rhythmic motion.

\textbf{Memory.} Memory is conceptualized differently in each framework. In GWT, memory is episodic, recruited into the workspace when contextually relevant \citep{dehaene2001towards}. In AIF, memory is statistical: priors embedded in hierarchical generative models constrain inference \citep{friston2017graphical}. RSVP--CPG, by contrast, models memory as proxy loops entrained to the cortical gait. These rhythmic subchains act like phonological loops, maintaining resonant traces of past trajectories and biasing future steps. Memory thus persists not through storage or priors, but through ongoing rhythmic resonance.

\textbf{Dynamics of Motion.} In GWT, dynamics are discrete: a representation either wins access or remains local. In AIF, dynamics are continuous but converge toward equilibrium: prediction errors are minimized until the system balances expectations with input \citep{friston2010free}. RSVP--CPG resists both discreteness and convergence. Its dynamics are oscillatory and non-equilibrium: cognition as perpetual descent along entropy gradients. Like gait, the process is sustained not by reaching stability but by repeatedly “falling forward” into the next phase.

\textbf{REM and Dreaming.} Dreaming provides a particularly sharp lens for distinguishing theoretical commitments. Global Workspace Theory interprets REM dreams as a workspace filled in the absence of sensory input: imagery generated internally gains ignition and broadcast, creating a stage populated by simulations rather than perception \citep{dehaene2014consciousness}. Active Inference instead frames REM as the generative model “free-running”: priors and predictions unfold without the anchoring of sensory evidence, enabling the exploration of hypothesis spaces otherwise suppressed in waking life \citep{hobson2012waking}.

The RSVP--CPG framework introduces a different emphasis. REM is not simply a workspace running offline, nor a generative model free-wheeling in imagination. Instead, REM is a shift in rhythm. Cortical cycles continue to turn, but their entrainment to CPG chains loosens. In waking and non-REM states, cortical cycles drive synchronized chains — linking memory loops, motion rhythms, and attentional steps into coherent gaits. In REM, these chains desynchronize: subloops drift out of phase, losing their tight alignment to the global cortical cycle. The result is a state in which multiple gaits coexist but fail to coordinate — a parade of out-of-step rhythms whose collisions manifest phenomenologically as the bizarre juxtapositions, discontinuities, and dreamlike logic of REM.

This rhythmic interpretation shifts the explanatory burden. In GWT and AIF, dream bizarreness is attributed to content — the workspace filling with internally generated images, or the generative model exploring implausible priors. In RSVP--CPG, bizarreness is an emergent property of form: it arises from phase misalignment in a rhythmic substrate. The dream is not “what happens when the lights stay on without input” (GWT) or “what happens when the model runs unconstrained” (AIF), but “what happens when the orchestra continues to play but the instruments slip out of sync.” This formal explanation makes new predictions: dream vividness should track cortical cycle strength, while dream bizarreness should track cross-chain phase jitter — variables now measurable in MEG-derived cycle metrics \citep{van2025cortical}.

Beyond REM: Explaining the Limits of GWT and AIF. The REM case study highlights more than a difference in explanatory flavor: it exposes the structural limits of both Global Workspace Theory and Active Inference. GWT offers a phenomenological description of conscious access, but its stage/spotlight metaphor lacks a mechanism for why contents fall out of sync in REM or why cognition unfolds rhythmically even when the “spotlight” is absent. AIF, meanwhile, provides a mathematically rigorous account of predictive regulation, yet its commitment to equilibrium means it struggles to account for paradoxical regimes where synchronization deliberately fails, such as REM sleep or certain dissociative states.



The RSVP--CPG framework fills this gap by providing the missing oscillatory skeleton. Cycles are not accidents layered atop a predictive engine or a broadcast stage; they are the rhythmic substrate that makes both possible. Broadcasting can be understood as a momentary alignment of chains that amplifies one content; prediction-error minimization can be reframed as the entropic pacing that constrains how far chains drift from synchrony. In this sense, RSVP--CPG does not merely complement GWT and AIF but explains why they work at all — and why they break down under conditions like REM.

\textbf{Synthesis.} Taken together, these three frameworks provide complementary perspectives. GWT offers a phenomenological account of conscious access. AIF provides the statistical grammar of predictive computation. RSVP--CPG supplies the oscillatory skeleton — the rhythmic substrate on which both broadcasting and prediction ride. By repositioning consciousness not as spotlight or equilibrium but as gait, RSVP--CPG reframes cognition as perpetual motion: an ordered fall through time, sustained by entropy itself.


\subsection{Entropy pacing and sleep regimes}

A slow RSVP pacing variable $S$ modulates synchrony in both cortical and CPG layers:
\begin{equation}
\dot S ;=; \varepsilon\Big(\beta_1 |z|^2 + \beta_2 \Gamma_\phi - \rho,(S-S_0)\Big),
\qquad 0<\varepsilon\ll1,
\label{eq:S}
\end{equation}
with $\Gamma_\phi$ defined in \eqref{eq:feedback} as CPG synchrony.

Coupling parameters are gated by $S$:
\begin{equation}
\eta(S) = \eta_0 + \chi_\eta S,\quad
\kappa(S) = \kappa_0 + \chi_\kappa S,\quad
\gamma(S) = \gamma_0 - \chi_\gamma S.
\label{eq:Scoupling}
\end{equation}

\paragraph{Wakefulness (synchronized regime).}
At moderate $S\approx S_0$, cortical and CPG couplings are strong
($\eta,\kappa$ large), $\gamma$ small, so cycles remain phase-locked.
This corresponds to normal cognition: cortical cycles gait motor and mnemonic CPGs in ordered sequence.

\paragraph{REM / paradoxical sleep (desynchronization regime).}
During REM, $S$ increases beyond threshold, weakening $\eta,\kappa$ while amplifying $\gamma$.
Cortical skew drives cycles forward, but inter-chain synchrony collapses:

Γ
𝜙
↓
,
chains drift semi-independently
.
Γ
ϕ
	​

↓,chains drift semi-independently.

This produces multiple overlapping but misaligned loops, manifesting as dreamlike associative leaps.
Lamphrons (tensions) relax without lamphrodyne closure, consistent with paradoxical imagery and replay-like bursts.

\paragraph{NREM (damped regime).}
At low $S$, $\mu<0$ and CPG coupling dominates, so cortical cycles damp into quiescence (slow-wave oscillations).
This corresponds to deep sleep with reduced cognitive gait.

\subsection{Observables and empirical alignment}

\begin{itemize}
\item \textbf{Cycle strength (TINDA)} $;;\leftrightarrow;;$ amplitude of $z$ and synchrony $\Gamma_\phi$.
\item \textbf{Cycle rate} $;;\leftrightarrow;;$ effective $\Omega=\omega-\beta\mu/\alpha$, heritable via $\mu,\omega$ parameters.
\item \textbf{Asymmetry (FO)} $;;\leftrightarrow;;$ skew parameter $\gamma(S)$, breaking detailed balance.
\item \textbf{Behavioral phase effects} $;;\leftrightarrow;;$ cortical cycle windows $W_i(\vartheta_i)$ aligning with CPG proxies \eqref{eq:proxy}.
\item \textbf{REM signatures} $;;\leftrightarrow;;$ reduced $\Gamma_\phi$, broadened phase distribution of $z$, intermittent bursts.
\end{itemize}

\section{Phase-Reduced Dynamics of RSVP--CPG}

\subsection{From amplitude--phase to phase-only dynamics}

Each cortical/CPG oscillator can be expressed in polar coordinates
\[
z_i(t) \;=\; r_i(t)\,e^{i\vartheta_i(t)}.
\]
In the weak-coupling limit, amplitudes $r_i$ saturate near
$r_i \approx \sqrt{\mu/\alpha}$ and are slaved to the phase.
The RSVP--CPG system reduces to a phase-only model:
\begin{equation}
\dot{\vartheta}_i \;=\; \omega_i
+ \frac{K(S)}{N} \sum_{j=1}^N \sin(\vartheta_j - \vartheta_i)
+ \gamma(S),
\label{eq:phase}
\end{equation}
where:
\begin{itemize}
\item $\omega_i$ are the natural frequencies (heritable variation).
\item $K(S) = \eta(S) + \kappa(S)$ is the effective synchronizing strength,
modulated by the slow entropy pacing variable $S$.
\item $\gamma(S)$ is a skew parameter that breaks detailed balance and drives
forward ``gait-like'' progression.
\end{itemize}

\subsection{Order parameter and synchrony}

Define the Kuramoto order parameter:
\[
R e^{i\Psi} \;=\; \frac{1}{N} \sum_{j=1}^N e^{i\vartheta_j}.
\]
Here $R \in [0,1]$ measures global phase coherence, while $\Psi$ is the mean
phase.  
Within RSVP--CPG:
\begin{itemize}
\item \textbf{Wake:} $R \approx 1$ (synchronized gait).
\item \textbf{REM:} $R \downarrow$, oscillators drift incoherently.
\item \textbf{NREM:} amplitudes $r_i \to 0$, global $R$ loses meaning.
\end{itemize}

\subsection{REM as bifurcation}

As $S$ increases, coupling $K(S)$ weakens.  
Kuramoto theory gives a critical threshold:
\begin{equation}
K(S_c) \;=\; \frac{2}{\pi g(0)},
\end{equation}
where $g(\omega)$ is the distribution of natural frequencies.  

\begin{itemize}
\item If $K(S) > K(S_c)$, oscillators synchronize ($R>0$).
\item If $K(S) < K(S_c)$, oscillators desynchronize ($R \approx 0$).
\end{itemize}

Thus, when entropy pacing $S$ crosses $S_c$, the system undergoes a
\emph{synchrony--incoherence bifurcation}, marking the onset of REM sleep.

\subsection{Phase diagram of RSVP sleep regimes}

The RSVP--CPG framework therefore predicts three dynamical regimes:
\begin{enumerate}
\item \textbf{NREM (low $S$):} amplitudes damp, $r_i \to 0$, cycles suppressed.
\item \textbf{Wake (intermediate $S$):} $K(S) > K_c$, $R\approx 1$,
yielding a locked cognitive gait.
\item \textbf{REM (high $S$):} $K(S) < K_c$, synchrony collapses,
$R \downarrow$, producing drifting, dreamlike associations.
\end{enumerate}

\subsection{Interpretation}

The phase-reduced model makes explicit that:
\begin{itemize}
\item Cognition corresponds to the synchronized regime of cortical/CPG
oscillators (stable gait).
\item Dreams emerge naturally as desynchronization (gait incoherence).
\item Individual variation is encoded in $g(\omega)$, consistent with
heritability of cycle asymmetry.
\end{itemize}

Thus REM sleep is not an epiphenomenon but a mathematically precise
\emph{bifurcation} in the RSVP oscillator network, driven by the slow
entropy variable $S$.

\section{Multiplex Kuramoto Formalization of RSVP--CPG}
\label{sec:multiplex-kuramoto}

We extend the phase-reduced RSVP--CPG model to a two-layer \emph{multiplex} Kuramoto network with explicit cortical and memory/CPG layers. The multiplex formulation makes explicit how entropy pacing $S$ modulates intra- and inter-layer couplings and how REM emerges from selective decoupling.

\subsection{Model definition}

Let the cortical layer consist of $N_c$ oscillators with phases $\{\vartheta_i^{(c)}\}_{i=1}^{N_c}$ and the memory/CPG layer consist of $N_m$ oscillators with phases $\{\vartheta_j^{(m)}\}_{j=1}^{N_m}$. The phase dynamics are
\begin{align}
\dot{\vartheta}_i^{(c)} &= \omega_i^{(c)} 
  + \frac{K_{cc}(S)}{N_c}\sum_{p=1}^{N_c} \sin\!\big(\vartheta_p^{(c)}-\vartheta_i^{(c)}-\delta_{ip}^{(cc)}\big)
  + \frac{K_{cm}(S)}{N_m}\sum_{q=1}^{N_m} \sin\!\big(\vartheta_q^{(m)}-\vartheta_i^{(c)}-\delta_{iq}^{(cm)}\big)
  + \gamma_c(S) + \sigma_c\,\xi_i^{(c)}(t),
\label{eq:multi-cortex}\\
\dot{\vartheta}_j^{(m)} &= \omega_j^{(m)}
  + \frac{K_{mm}(S)}{N_m}\sum_{q=1}^{N_m} \sin\!\big(\vartheta_q^{(m)}-\vartheta_j^{(m)}-\delta_{jq}^{(mm)}\big)
  + \frac{K_{mc}(S)}{N_c}\sum_{p=1}^{N_c} \sin\!\big(\vartheta_p^{(c)}-\vartheta_j^{(m)}-\delta_{jp}^{(mc)}\big)
  + \gamma_m(S) + \sigma_m\,\xi_j^{(m)}(t).
\label{eq:multi-memory}
\end{align}
Here:
\begin{itemize}
  \item $\omega^{(\cdot)}$ are natural frequencies (drawn from layer-specific distributions $g_c(\omega),g_m(\omega)$).
  \item $K_{\alpha\beta}(S)$ are entropy-paced coupling strengths ($\alpha,\beta\in\{c,m\}$).
  \item Phase-lag matrices $\delta_{ab}^{(\cdot)}$ allow chain-specific lags and asymmetric hand-off.
  \item $\gamma_{c,m}(S)$ are layer skew/drift terms (forward drive).
  \item $\sigma_{c,m}\,\xi(t)$ are independent white-noise terms modelling spontaneous bursts (e.g. hippocampal noise); $\langle\xi(t)\rangle=0,\ \langle\xi(t)\xi(t')\rangle=\delta(t-t')$.
\end{itemize}

\subsection{Entropy pacing \& REM decoupling}

Entropy pacing is governed by a slow variable $S(t)$ (cf. main text). We model its effect on couplings as smooth monotone maps:
\begin{equation}
K_{\alpha\beta}(S) = K_{\alpha\beta}^0 \, f_{\alpha\beta}(S),\qquad
f_{\alpha\beta}:\mathbb{R}\to\mathbb{R}_+, \quad f_{\alpha\beta}(S_0)=1,
\end{equation}
with typical choices $f_{\alpha\beta}(S)=\exp(-\chi_{\alpha\beta} (S-S_0))$ or sigmoids. REM-like decoupling is implemented by making inter-layer couplings particularly sensitive to $S$:
\[
f_{cm}(S), f_{mc}(S) \downarrow \quad \text{rapidly as } S \uparrow,
\]
so that as $S$ crosses a threshold $S_c$ the effective cortex-to-memory drive collapses while intra-layer cortical coupling $K_{cc}(S)$ may remain moderate (yielding persistent cortical cycling with reduced cross-layer alignment).

\subsection{Layer order parameters and readouts}

Define complex order parameters for each layer:
\begin{align}
R_c e^{i\Psi_c} &= \frac{1}{N_c}\sum_{p=1}^{N_c} e^{i\vartheta_p^{(c)}},\qquad
R_m e^{i\Psi_m} = \frac{1}{N_m}\sum_{q=1}^{N_m} e^{i\vartheta_q^{(m)}}.
\end{align}
Interpretation:
\begin{itemize}
  \item $R_c,R_m\in[0,1]$ measure intra-layer synchrony (gait coherence).
  \item Phase difference $\Delta\Psi=\Psi_c-\Psi_m$ indexes cross-layer alignment (binding of cortical windows to CPG proxies).
  \item Cross-layer coherence can also be quantified as
  \[
  C_{cm} = \left|\frac{1}{N_c N_m}\sum_{p,q} e^{i(\vartheta_p^{(c)}-\vartheta_q^{(m)})}\right|,
  \]
  which decreases in REM-like decoupling.
\end{itemize}

\subsection{Reduced mean-field description (weak heterogeneity)}

Under weak heterogeneity and mean-field coupling, one can approximate the cortical layer by the Ott–Antonsen (OA) ansatz (similarly for memory layer) to obtain closed low-dimensional dynamics for $(R_c,\Psi_c)$ and $(R_m,\Psi_m)$. The OA reduction yields schematic equations (omitting derivation):
\begin{align}
\dot R_c &= -\Delta_c R_c + \tfrac{K_{cc}(S)}{2} R_c (1-R_c^2) + \tfrac{K_{cm}(S)}{2}\,F_{cm}(R_m,R_c,\Delta\Psi) + \mathcal{O}(\sigma_c^2),\\
\dot \Psi_c &= \bar\omega_c + \gamma_c(S) + \tfrac{K_{cm}(S)}{2}\,G_{cm}(R_m,R_c,\Delta\Psi) + \mathcal{O}(\sigma_c^2),
\end{align}
with analogous equations for $(R_m,\Psi_m)$. Here $\Delta_{c,m}$ denote frequency dispersion widths; $F_{cm},G_{cm}$ are geometric coupling functions depending on phase difference. The OA reduction makes it straightforward to analyze fixed points and bifurcations as $S$ varies.

\subsection{Dynamical regimes and signatures}

Using the full model \eqref{eq:multi-cortex}--\eqref{eq:multi-memory} or the OA reduction, RSVP predicts:
\begin{itemize}
  \item \textbf{Wake (low--intermediate $S$):} $K_{cc},K_{mm},K_{cm},K_{mc}$ large $\Rightarrow$ $R_c\approx 1,R_m\approx 1$, $C_{cm}$ high (locked gait; memories carried reliably).
  \item \textbf{REM (high $S$ beyond $S_c$):} rapid drop of $K_{cm},K_{mc}$ $\Rightarrow$ $R_c$ may remain moderate (cortical cycles persist) while $R_m$ and $C_{cm}$ fall or fragment; chimera states (coexistence of coherent and incoherent subnetworks) and increased phase diffusion emerge, consistent with dream-like recombination.
  \item \textbf{NREM (very low $S$ or large damping):} effective damping of amplitudes (outside phase model regime) leads to suppressed cycling; $R_c,R_m$ collapse toward incoherent/noisy baselines.
\end{itemize}

\subsection{Noise, chimera states, and partial synchrony}

Noise terms $\sigma_{c,m}\xi$ play multiple roles:
\begin{enumerate}
  \item In REM, with weakened inter-layer coupling, noise drives phase wandering and stochastic recombination of CPG proxies (explains bizarre, randomly associated dream content).
  \item Intermediate regimes admit chimera states where subsets synchronize ($R_{\text{sub}}>0$) while others desynchronize; these are plausible mechanistic accounts of partial or fragmented conscious content (dreams with coherent strands).
  \item Noise-induced switching between metastable synchronized patterns can model trial-to-trial variability in cycle ordering observed empirically.
\end{enumerate}

\subsection{Observables and empirical mapping}

The multiplex model yields directly testable predictions for MEG/EEG/TINDA-style analyses:
\begin{itemize}
  \item \textbf{FO asymmetry matrices:} skew terms $\gamma_{c,m}(S)$ and phase-lag structure $\delta$ map to preferred transition directions; estimated FO asymmetries should scale with $\gamma_c$ and inter-layer topology.
  \item \textbf{Cycle strength and rate:} cortical order parameter $R_c$ amplitude and mean phase velocity $\dot\Psi_c$ map to empirically measured cycle strength and rate; heritability of $\omega$ distributions predicts trait-like rate stability.
  \item \textbf{Cross-layer coherence drop in REM:} $C_{cm}$ should sharply decrease at REM onset while cortical intra-layer metrics (e.g. $R_c$) may remain non-zero.
  \item \textbf{Chimera signatures:} spatially localized coherence (high $R$ in subnetworks, low globally) is detectable as partial-state synchronization in HMM/TINDA analyses.
\end{itemize}

\subsection{Simulation and fitting notes}

\begin{itemize}
  \item Use Euler–Maruyama integration with timestep $\Delta t \ll \min\{K_{\alpha\beta}^{-1},\sigma^{-2}\}$.
  \item Fit layer-level parameters $(K^0_{\alpha\beta},\chi_{\alpha\beta},\gamma_{c,m}^0,\Delta_{c,m})$ by minimizing mismatch between model and empirical statistics: cycle rate distribution, FO asymmetry matrix, $R_c(t)$ histograms, and cross-layer coherence $C_{cm}$.
  \item Bayesian hierarchical fitting can capture subject-level $\omega$ draws and heritable components by including random effects on $\omega$ or on a latent $S_*$.
\end{itemize}

\subsection{Summary}

The multiplex Kuramoto formalization captures the key RSVP mechanistic claims: cortical cycles function as a synchronized gait in wake, memory/CPG chains are coupled proxies that ride cortical windows, and REM emerges when entropy pacing selectively decouples inter-layer interactions. The framework admits rich dynamical phenomena (chimera states, noise-driven recombination, ultradian cycling when $S$ is itself dynamical) and maps cleanly to MEG/EEG observables.

To broaden this comparison, we engage with additional theories in the philosophy of mind and neuroscience.

\subsection{Broader Theoretical Landscape}

Integrated Information Theory (IIT) posits that consciousness arises from integrated information \citep{tononi2004information}. In RSVP--CPG, the rhythmic integration of cycles provides a dynamic basis for IIT's Φ measure.

Orchestrated Objective Reduction (Orch-OR) suggests quantum processes in microtubules contribute to consciousness \citep{hameroff2014consciousness}. RSVP--CPG could incorporate quantum oscillations in CPG dynamics.

Attention Schema Theory views consciousness as an internal model of attention \citep{webb2015attention}. The gait metaphor extends this to rhythmic modeling of self and other.

The Dynamic Core Hypothesis emphasizes reentrant neural interactions \citep{tononi1998consciousness}. RSVP--CPG's chained oscillators provide a rhythmic implementation of this core.

\subsection{Critical Engagement with Existing Frameworks}

Potential criticisms of RSVP--CPG include its reliance on metaphorical language. However, the mathematical formalization addresses this by providing testable predictions.

Limitations include the current focus on cortical levels; future work should integrate subcortical and peripheral rhythms.

\section{RSVP Interpretation of Cortical Cycles}

Recent magnetoencephalography studies have revealed that large-scale cortical networks activate in robust cyclical sequences, with asymmetric transitions, stable ordering, and strong heritability of cycle rate \citep{van2025cortical}. These findings call for a theoretical framework that explains not only the existence of such cycles but also their role in cognition, memory, and sleep.

1. Cycles as Entropic Smoothing (Lamphron–Lamphrodyne Dynamics). In RSVP theory, lamphron–lamphrodyne dynamics describe how local energy differentials are smoothed into coherent field patterns through recursive entropic redistribution. The Oxford study shows that cortical networks transition in asymmetric but ordered sequences, forming robust cycles. This mirrors RSVP’s idea that network activity is not stochastic, but guided by an entropic descent that smooths and re-integrates activity into temporally ordered attractors.

Lamphron: The momentary tension or energy differential between networks.

Lamphrodyne: The relaxation process that pushes activity into the next viable network in the cycle.

Thus, cortical cycles can be seen as a mesoscale instantiation of lamphron–lamphrodyne smoothing: the brain keeps activity “falling forward” along constrained entropic pathways.

2. Φ–𝒗–S Attractor Oscillations. In RSVP’s field model: Φ (scalar density): encodes baseline excitability or readiness of cortical territories. 𝒗 (vector flow): captures directed transitions of activation between networks. S (entropy field): tracks uncertainty and redundancy in network coordination.

The Oxford findings suggest that cortical networks follow Φ–𝒗–S cycles, where: Φ rises locally (network primed for activation). 𝒗 drives the transition asymmetrically toward a preferred successor. S regulates the overall pacing, ensuring the cycle doesn’t collapse into chaos or stall.

In this view, the cyclical pattern is not just clock-like but a Lyapunov-stable attractor orbit in RSVP’s field space.

3. Recursive Causality and Temporal Windows. RSVP emphasizes recursive causality—future constraints feed back into present organization. The Oxford group observed that cycles operate on timescales faster than cognition but still constrain behavioral order.

RSVP interpretation: cycles act as temporal scaffolds, partitioning cognitive operations into windows where memory, perception, or movement are most efficiently executed.

This makes cycles equivalent to RSVP’s recursive tilings: dynamic partitions of field space that shape when certain computations can occur.

4. Clinical and AI Implications. Clinical: Disorders like Alzheimer’s or ADHD may reflect cycle disruptions, where the entropic smoothing fails (lamphron accumulates without lamphrodyne relaxation), leading to cognitive bottlenecks or chaotic transitions.

AI \& RSVP-AI: In RSVP-based AI, cyclical Φ–𝒗–S attractors could serve as intrinsic schedulers, preventing overfitting by enforcing entropic pacing of inference. This suggests cortical cycles are nature’s own implementation of RSVP’s semantic merge operators, ensuring coherence across distributed modules.

5. Philosophical Perspective. RSVP treats cognition as an entropic computation embedded in physical fields. The discovery of cortical cycles validates the idea that constraint relaxation, not expansion, is the organizing principle of intelligence. Just as in cosmology RSVP denies universal expansion and instead models constant entropic smoothing, in neuroscience RSVP reframes cognition as structured recurrence within fixed capacity, not linear accumulation of “processing power.”

To strengthen this interpretation, we provide a detailed stability analysis and information-theoretic foundations.
\subsection{Stability Analysis of Limit Cycles}

For RSVP--CPG to serve as a robust model of cortical dynamics, it must be shown that its cyclical solutions are mathematically stable. We approach this through Lyapunov analysis and bifurcation theory.

Consider the cortical cycle dynamics in equation (\ref{eq:cortex}), a complex Hopf network with skew-symmetric coupling. Linearization around the origin reveals that for $\mu > 0$ and $\alpha > 0$, the trivial fixed point is unstable, and trajectories are repelled outward. Nonlinear cubic terms then constrain growth, producing a bounded oscillation: a limit cycle. This is the classical supercritical Hopf bifurcation.

To establish stability, we define a Lyapunov function
\begin{equation}
V(z) = \frac{1}{2}\|z\|^2,
\end{equation}
whose derivative along trajectories yields
\begin{equation}
\dot V = \Re\langle z, \dot z \rangle
= \mu \|z\|^2 - \alpha \|z\|^4 + \ldots
\end{equation}
For $\mu > 0$ and $\alpha > 0$, $\dot V > 0$ when $\|z\|$ is small and $\dot V < 0$ when $\|z\|$ is large, implying that trajectories are repelled from the origin but attracted to a finite-radius cycle. This confirms global asymptotic stability of the limit cycle.

Parameter variation induces bifurcations that map naturally onto cognitive regimes. Increasing $\gamma$ strengthens asymmetry, enlarging cycle strength and ordering. Increasing $\eta$ damps transitions, flattening cycles (as in NREM). Decreasing $\kappa$ in the CPG layer destabilizes cross-chain synchrony, producing desynchronized regimes (as in REM). These bifurcations correspond to experimentally observed transitions across wake, non-REM, and REM sleep.

Thus, RSVP--CPG provides not just a metaphor but a mathematically rigorous account of stable cyclic dynamics in the brain. The existence and stability of its limit cycles ensure that the cortical gait persists under perturbation — an essential property for both cognition and survival.


\subsection{Information-Theoretic Foundations}

The RSVP--CPG framework treats cortical cycles as processes of entropic smoothing: rhythmic progressions that reduce local unpredictability while preserving global variability. To render this precise, we must situate the model within an information-theoretic framework.

Consider a cortical network state $X_t$ at time $t$, with distribution $p(x_t)$. The entropy of this distribution,
\begin{equation}
H(X_t) = - \sum_x p(x_t) \log p(x_t),
\end{equation}
quantifies the uncertainty of the network. RSVP dynamics proceed by coupling successive states such that the entropy of local subsystems decreases while the joint entropy across cycles remains high. In other words, each cortical cycle functions as a step in an information pipeline: smoothing within-cycle uncertainty while maintaining between-cycle richness.

Formally, we can describe this as minimizing the conditional entropy
\begin{equation}
H(X_{t+1} \,|\, X_t),
\end{equation}
subject to a constraint on the global entropy rate
\begin{equation}
\lim_{T \to \infty} \frac{1}{T} H(X_1, \dots, X_T) \geq \kappa,
\end{equation}
where $\kappa$ ensures diversity across cycles. This balance resembles the free energy principle \citep{friston2010free}, but with a distinctive twist: RSVP--CPG dynamics are not convergent equilibria but rhythmic limit cycles. Rather than minimizing surprise outright, they channel entropy into oscillatory gaits, ensuring continual renewal.

This formalism explains the dual character of cortical cycles: locally predictable, globally variable. It also provides testable predictions: disruptions of entropy balance (e.g., excess within-cycle noise or excessive global regularity) should correspond to pathological states such as epilepsy (hyper-synchrony) or schizophrenia (hypo-synchrony). Thus, RSVP--CPG casts cortical oscillations as information-theoretic engines, cycling between entropy minimization and entropy preservation to sustain the flow of cognition.


Entropic smoothing is formalized as minimization of a specified entropy functional, connecting to variational principles in physics.

\section{An RSVP--CPG Unified Model of Cognitive Gait, Memory Loops, and REM Desynchronization}

Empirical anchor. Large-scale cortical functional networks activate in robust cycles (periods $\sim$300--1{,}000\,ms), with asymmetric transition probabilities, reproducible ordering across datasets, behavioral phase relevance, and strong heritability of cycle rate. These properties are quantified via TINDA-derived FO-asymmetry matrices, cycle strength, and cycle rate \citep{van2025cortical}.

\subsection{State variables and scales}
We model three coupled layers: (i) a cortical \emph{cycle layer} $z\in\mathbb{C}^K$ capturing $K$ canonical large-scale networks; (ii) a \emph{CPG layer} of $M$ oscillators $\phi\in\mathbb{S}^1$ arranged in chains (motor and mnemonic); (iii) a slow \emph{entropy/pacing} variable $S$ regulating synchrony. RSVP fields $(\Phi,\mathbf{v},S)$ are represented in reduced form: $\Phi\leftrightarrow |z|$ (excitability), $\mathbf{v}\leftrightarrow \nabla\arg z$ (directed flow), and $S$ as a slow control of gains and couplings.

\subsection{Cortical cycle layer: complex Hopf network with directed skew}
Let $G=(V,E)$, $|V|=K$, with Laplacian $L$ and skew-symmetric $K=-K^\top$ encoding preferential transitions (broken detailed balance). The cortical dynamics are
\begin{equation}
\dot z \;=\; (\mu + i\omega)\,z \;-\; (\alpha + i\beta)\,(z\odot\bar z)\odot z
\;-\; \eta\,L z \;+\; i\,\gamma\,K z \;+\; \rho_z\,U_z(t),
\label{eq:cortex}
\end{equation}
with $z\in\mathbb{C}^K$, elementwise cubic nonlinearity, diffusive coherence $-\eta L z$, and a \emph{skew rotator} $i\gamma K z$ generating the global cycle and asymmetric state order. For $\mu>0$ and $\alpha>0$, \eqref{eq:cortex} admits a stable limit cycle with mean period $T\approx 2\pi/\Omega$, $\Omega=\omega-\beta\,\mu/\alpha$, matching the 300--1{,}000\,ms band. The parameter $\gamma>0$ controls FO-asymmetry magnitude and thereby cycle strength. $U_z$ captures exogenous/task locking.

\subsection{CPG chains: coupled-phase oscillators for motor \emph{and} mnemonic loops}
We represent chained CPGs (motor or mnemonic/phonological) as rings/paths of $M$ phases $\phi=(\phi_1,\dots,\phi_M)$:
\begin{equation}
\dot\phi_j \;=\; \omega_j \;+\; \frac{\kappa}{d_j}\!\sum_{k\in\mathcal{N}(j)}\!\sin(\phi_k-\phi_j-\delta_{jk})
\;+\; \lambda\,\Xi_j(z) \;+\; \rho_\phi\,U_\phi(t) \;+\; \sigma_\phi\,\xi_j(t).
\label{eq:cpg}
\end{equation}
Here $\omega_j$ are natural frequencies; nearest-neighbor (or graph) coupling $\kappa>0$ induces gait-like phase relations; $\delta_{jk}$ encodes chain-specific phase lags; $\Xi_j(z)$ is a cortical drive from \eqref{eq:cortex} (defined below); $U_\phi$ external inputs; $\xi_j$ noise.

\paragraph{Mnemonic phonological loop.} A phonological/memetic item $m$ is maintained as a phase-coded proxy bound to a CPG subchain $\mathcal{C}\subset\{1,\dots,M\}$:
\begin{equation}
\textstyle \Pi_m(t) \;=\; \frac{1}{|\mathcal{C}|}\sum_{j\in\mathcal{C}} \cos\big(\phi_j(t)-\theta_m\big),
\quad \theta_m\in[0,2\pi)
\label{eq:proxy}
\end{equation}
with $\Pi_m$ reactivated when cortical phases enter a receptive \emph{cycle window} (defined from $\arg z$ below). Sequencing arises from ordered windows across the cortical cycle.

\subsection{Cross-layer coupling (RSVP binding)}
We bind cognitive \emph{cycle phase} to CPG chains via two couplings:

\paragraph{(A) Cortex $\rightarrow$ CPG (driving).} Let $\vartheta_i=\arg z_i$ and define $W_i(\vartheta)$ as narrow phase windows marking the preferred activation of network $i$ along the global cycle. Map these to CPG targets via an assignment $\mathcal{M}\!:\, i\mapsto \mathcal{M}(i)\subset\{1,\dots,M\}$. Then
\begin{equation}
\Xi_j(z) \;=\; \sum_{i=1}^K g_{ij}\, W_i(\vartheta_i), \qquad
g_{ij} =
\begin{cases}
g_+ & j\in \mathcal{M}(i)\\
0 & \text{otherwise.}
\end{cases}
\label{eq:drive}
\end{equation}
During wake, $g_+$ aligns CPG phases to cortical cycle windows, yielding gaited cognition.

\paragraph{(B) CPG $\rightarrow$ cortex (feedback).} CPG coherence feeds back as a gain on cortical amplitude and skew:
\begin{equation}
\mu \;=\; \mu_0 + \chi_\mu\,\Gamma_\phi,\qquad
\gamma \;=\; \gamma_0 + \chi_\gamma\,\Gamma_\phi,\qquad
\Gamma_\phi \;=\; \frac{1}{M}\sum_{j}\Big|\frac{1}{d_j}\sum_{k\in\mathcal{N}(j)} e^{i(\phi_k-\phi_j)}\Big|.
\label{eq:feedback}
\end{equation}
Higher CPG synchrony $\Gamma_\phi$ strengthens cycle rate/strength, consistent with faster, more ordered cycles linked to behavior.

\subsection{Entropy pacing and sleep regimes}
A slow RSVP pacing variable $S$ modulates the coherence parameters:
\begin{equation}
\dot S \;=\; \varepsilon\big(\beta_1\|z\|^2 + \beta_2 \Gamma_\phi - \rho(S-S_0)\big),\quad
\eta = \eta(S),\ \kappa=\kappa(S),\ \gamma=\gamma(S),\ \sigma_\phi=\sigma_\phi(S),\ \rho_z=\rho_z(S).
\label{eq:pacing}
\end{equation}
\textbf{Wake (synchronized chains).} $S$ sets high $\kappa$, moderate $\eta$, elevated $\gamma$: cortical ordering drives CPG chains; mnemonic proxies ride stable loops; cycle phase predicts reaction time and replay windows.

\textbf{NREM (overdamped).} $S$ increases $\eta$, reduces $\gamma$ and $\rho_z$: cycle amplitude damps; CPG chains remain largely synchronized but with lowered drive.

\textbf{REM/paradoxical (desynchronization).} $S$ reduces $\kappa$ and increases $\sigma_\phi$, weakens $g_+$ and feedback \eqref{eq:feedback}; cortical limit cycle persists but cross-layer alignment loosens. Multiple subchains drift in and out of phase, producing associative recombination and paradoxical imagery while preserving fast cortical cycling. This realizes the \emph{desynchronized CPG chains} hypothesis for REM within RSVP.

\subsection{Measurement model and empirical mapping}
\paragraph{MEG/TINDA observables.} From $z(t)$, define state activations by thresholding $|z_i|$ and estimate FO-asymmetry matrices $A_{mn}$ over variable interstate intervals (TINDA). The model predicts (i) significant global cycle strength from the joint action of $i\gamma K$ and $-\eta L$, (ii) strongest asymmetries at longer intervals, and (iii) subject-specific cycle rate/strength that correlate with covariates and behavior \citep{van2025cortical}.

\paragraph{Behavioral coupling.} Let $\mathcal{I}_{\text{vis}}$ mark high-power visual/attentional states and $\mathcal{I}_{\text{low}}$ low-power sensorimotor states on the cycle. Define a phase score $\psi(t)=\sum_{i\in\mathcal{I}_{\text{vis}}}\!W_i(\vartheta_i(t))- \sum_{i\in\mathcal{I}_{\text{low}}}\!W_i(\vartheta_i(t))$. Model log-RT as $\log \mathrm{RT} \sim \mathcal{N}(\alpha - \beta\,\psi(t-500\,\mathrm{ms}),\sigma^2)$, reproducing the reported opposite sign associations by cycle phase \citep{van2025cortical}.

\subsection{Predictions and tests}
\begin{enumerate}
\item \textbf{REM test.} In REM, $\kappa\!\downarrow$, $\sigma_\phi\!\uparrow$, $g_+\!\downarrow$: cycle strength (cortex) remains nonzero but \emph{cross-layer} synchrony falls; mnemonic proxies $\Pi_m$ show increased phase jitter across loops, with dream-like recombinations.
\item \textbf{Replay locking.} Memory replay events cluster in the ``north'' quadrant (DMN/alpha) cycle phases and are anti-phased with low-power sensorimotor states \citep{van2025cortical}.
\item \textbf{Heritability.} Individual cycle rate correlates with baseline $\omega$ and with steady-state $S_*$; a Bayesian ACE model on $\log$ rate yields high $h^2$, replicating twin results \citep{van2025cortical}.
\item \textbf{Intervention.} Phase-locked stimulation targeting $W_i(\vartheta_i)$ advances/delays the cycle; transient $\kappa$ boosting re-synchronizes chains (therapeutic for REM dysregulation).
\end{enumerate}

\subsection{Numerical scheme (graph form)}
A split-step Euler integrator for \eqref{eq:cortex}--\eqref{eq:cpg}--\eqref{eq:pacing}:
\[
\begin{aligned}
z^{t+\Delta t} &= z^t + \Delta t\big[(\mu+i\omega) z^t - (\alpha+i\beta)(z^t\!\odot\!\bar z^{\,t})\!\odot\! z^t
- \eta L z^t + i\gamma K z^t + \rho_z U_z^t \big],\\
\phi_j^{t+\Delta t} &= \phi_j^t + \Delta t\Big[\omega_j + \frac{\kappa}{d_j}\!\sum_{k\in\mathcal{N}(j)}\!
\sin(\phi_k^t-\phi_j^t-\delta_{jk}) + \lambda\,\Xi_j(z^t) + \rho_\phi U_{\phi,j}^t \Big] + \sqrt{\Delta t}\,\sigma_\phi \xi_j^t,\\
S^{t+\Delta t} &= S^t + \Delta t\big[\varepsilon(\beta_1\|z^t\|^2 + \beta_2 \Gamma_\phi^t - \rho(S^t-S_0))\big],
\end{aligned}
\]
with adaptive $\Delta t$ constrained by maximal cubic and coupling gains.

\paragraph{Bayesian calibration.} Estimate $(\Omega,\gamma,\eta)$ from FO-asymmetry matrices and cycle strength; recover $S$ by regressing cycle rate/strength and cross-layer coherence; fit ACE on $\log$ rate to quantify $h^2$, mirroring the empirical pipeline \citep{van2025cortical}.

To expand the model, we include stochastic terms and advanced stability analysis.
\subsection{Stochastic Analysis}

Neural oscillations are never perfectly periodic; they are embedded in noisy biological substrates. Ion channel fluctuations, synaptic release variability, and global state changes all inject stochasticity into cortical rhythms. For any model of cyclical cognition to be biologically plausible, it must capture not just deterministic limit cycles but also their perturbation by noise. This motivates a stochastic formulation of the RSVP--CPG equations.

Formally, we recast the cortical cycle dynamics (\ref{eq:cortex}) and CPG chain equations (\ref{eq:cpg}) as stochastic differential equations (SDEs). In It\^o form, the cortical variables evolve as
\begin{equation}
dz = f(z,t)\,dt + \Sigma_z\,dW_t,
\end{equation}
where $f(z,t)$ is the deterministic vector field defined previously, $\Sigma_z$ encodes noise covariance across networks, and $dW_t$ is a Wiener process. Likewise, CPG phases $\phi_j$ evolve under
\begin{equation}
d\phi_j = g_j(\phi,t)\,dt + \sigma_\phi\,dW_t^j,
\end{equation}
where $\sigma_\phi$ scales phase diffusion. The inclusion of noise transforms stable limit cycles into stochastic limit cycles, producing phase jitter and amplitude fluctuations that mirror the variability observed in MEG and EEG recordings.

Within the RSVP--CPG framework, noise has functional as well as disruptive roles. During wake, moderate noise allows exploration of near-by trajectories, enhancing cognitive flexibility. In non-REM sleep, increased damping suppresses noise, leading to more regular cycles. In REM sleep, by contrast, cross-chain couplings weaken while $\sigma_\phi$ rises, leading to desynchronization: chains drift in and out of alignment, producing the surreal recombinations characteristic of dreams. Thus, noise is not a bug but a mechanism for generating the diversity of cognitive states.

This stochastic framing makes concrete predictions. The variance of cycle strength should scale with $\sigma_\phi$ in a measurable way, predicting higher trial-to-trial variability of network activation sequences during REM than wake. Phase diffusion coefficients estimated from MEG could serve as biomarkers for cycle health, with disorders such as ADHD or schizophrenia reflecting excessive phase noise. Computationally, RSVP--CPG predicts that adding Gaussian noise to CPG chains should replicate both dream-like imagery and cognitive lapses in task performance.

By incorporating stochastic analysis, RSVP--CPG bridges deterministic dynamical models with the noisy reality of brain function. Noise becomes an integral part of the gait of thought, shaping when and how cycles drift, desynchronize, and reconfigure.

\subsection{Full Field Equations}

The RSVP--CPG framework formalizes cognition in terms of three interdependent fields: a scalar density field $\Phi(\mathbf{x},t)$, a vector flow field $\mathbf{v}(\mathbf{x},t)$, and an entropy field $S(\mathbf{x},t)$. Together these constitute a coupled dynamical system that governs the evolution of cortical cycles.

\paragraph{Scalar Field (Potential Density).}  
The scalar field $\Phi$ represents local excitation density, analogous to a potential surface. Its dynamics are given by:
\begin{equation}
\frac{\partial \Phi}{\partial t} = D_\Phi \nabla^2 \Phi - \lambda \Phi^3 + \mu \Phi + \alpha \, \nabla \cdot \mathbf{v} - \beta S,
\end{equation}
where $D_\Phi$ controls spatial diffusion, $\lambda$ enforces saturation, $\mu$ drives oscillatory onset, $\alpha$ couples scalar density to vector divergence, and $\beta$ implements entropy damping. This equation ensures that $\Phi$ evolves as a nonlinear oscillator constrained by flow and entropy.

\paragraph{Vector Field (Oscillatory Flow).}  
The vector field $\mathbf{v}$ represents rhythmic propagation of cortical activity, modeled as:
\begin{equation}
\frac{\partial \mathbf{v}}{\partial t} = D_v \nabla^2 \mathbf{v} - \gamma \|\mathbf{v}\|^2 \mathbf{v} + \kappa \nabla \Phi - \eta \nabla S + \mathbf{\Omega} \times \mathbf{v},
\end{equation}
where $D_v$ governs diffusion, $\gamma$ provides nonlinear damping, $\kappa$ aligns flow with scalar gradients, $\eta$ couples entropy gradients, and $\mathbf{\Omega}$ encodes torsional rotation (modeling oscillatory cycles). The cross product term ensures sustained rhythmicity, producing the gait-like oscillations characteristic of RSVP.

\paragraph{Entropy Field (Constraint Surface).}  
The entropy field $S$ captures information content and cycle constraint. Its dynamics are given by:
\begin{equation}
\frac{\partial S}{\partial t} = D_S \nabla^2 S + \theta \Phi^2 - \zeta \nabla \cdot \mathbf{v} - \rho S,
\end{equation}
where $D_S$ mediates diffusion, $\theta$ increases entropy with local activity, $\zeta$ reduces entropy through organized flow, and $\rho$ provides relaxation. The entropy field thus balances disorder and organization, mediating between scalar density and vector flow.

\paragraph{Coupled System.}  
Taken together, the RSVP--CPG system is a triadic coupling:
\begin{align}
\frac{\partial \Phi}{\partial t} &= f_\Phi(\Phi, \mathbf{v}, S), \\
\frac{\partial \mathbf{v}}{\partial t} &= f_v(\Phi, \mathbf{v}, S), \\
\frac{\partial S}{\partial t} &= f_S(\Phi, \mathbf{v}, S),
\end{align}
where each field simultaneously constrains and is constrained by the others. This structure guarantees oscillatory solutions: $\Phi$ oscillates through nonlinear potential wells, $\mathbf{v}$ cycles via torsion and divergence, and $S$ modulates coherence through entropic smoothing. The resulting trajectories form stable limit cycles—neural gaits—that embody cognition as rhythmic process.

\paragraph{Interpretation.}  
In phenomenological terms, $\Phi$ corresponds to local ignition (the intensity of representation), $\mathbf{v}$ to the directionality of thought (the vector of action or attention), and $S$ to cognitive context (the horizon of uncertainty). Cognition arises not from any single field, but from their cyclical interplay: scalar ignition initiates, vector flow propagates, entropy constraints regulate, and the cycle renews. This threefold dynamical structure provides the mathematical skeleton of RSVP--CPG.


\section{Bayesian Analysis for RSVP-Framed Cortical Cycles}

We present a Bayesian generative model that treats the Oxford MEG result as data, linking it to the RSVP limit-cycle model. (Detailed code and parameters omitted for brevity; see supplemental for Stan/NumPyro implementation.)
\section{Empirical Testability: Protocols and Predictions}
\label{sec:empirical}

This section specifies falsifiable tests of RSVP--CPG, with power analyses, acquisition parameters, preprocessing, statistical plans, controls, and preregistration items. Primary outcomes are (i) \emph{cycle strength} (global asymmetry magnitude), (ii) \emph{cycle rate} (Hz), (iii) \emph{phase jitter} (circular variance), (iv) \emph{cross-layer synchrony} (cortex{\textrightarrow}CPG coherence), and (v) \emph{behavioral phase locking} (RT/accuracy modulation by cycle phase).

\subsection{Study 1: Resting and Task MEG --- Detection and Reliability of Cycles}
\textbf{Participants.} $N{=}60$ healthy adults (18--35; balanced gender), right-handed, normal/corrected vision. Exclusions: neurological/psychiatric disorders, psychoactive meds.

\textbf{Acquisition.} Elekta 306-channel MEG; 1\,kHz; online 0.1--330\,Hz; Head Position Indicator (HPI) continuous tracking; 10\,min eyes-open rest (\emph{fixation}); 10\,min eyes-closed rest; 20\,min visual oddball task (ISI 1--1.5\,s; targets 20\%). T1 MRI (1\,mm) for source modeling.

\textbf{Preprocessing.} MaxFilter/SSS; bad channel interpolation; line-noise notch; ICA/SSP to remove blinks/ECG; bandpass 1--100\,Hz; source projection with LCMV beamformer to atlas ROIs (e.g., Schaefer-200).

\textbf{Cycle inference.} Hidden Markov Model (HMM) on ROI envelopes; compute transition matrices at variable lags (50--1500\,ms); derive FO-asymmetry; extract cycle decomposition; compute cycle strength and rate. Cross-validate (10-fold) within-subject; test--retest on a second session ($N{=}30$).

\textbf{Stats.} One-sample tests vs 0 for cycle strength; circular statistics for rate distributions; intraclass correlation (ICC) for reliability. FDR $q{<}.05$ across ROIs/lags.

\textbf{Power.} Pilot/previous effect sizes ($d{\sim}.45$) imply $N{=}60$ gives $1{-}\beta{>}0.90$ (two-tailed $\alpha{=}.05$) for cycle strength.

\textbf{Predictions.} Robust nonzero cycle strength; rates in the 0.8--3.5\,Hz band; higher asymmetry at longer lags; good reliability (ICC$>.6$).

\textbf{Controls.} Phase-randomized surrogates; HMM state-count sensitivity; head-movement regressors; vigilance covariates (pupil, alpha power).

\textbf{Preregistration.} OSF: primary outcomes, alpha control, exclusion rules, preprocessing pipeline.

\subsection{Study 2: Behavioral Coupling --- Phase Predicts Performance}
\textbf{Participants.} Subset of Study 1 ($N{=}50$).

\textbf{Tasks.} (a) Visual oddball; (b) 2-back WM; (c) Posner cueing. Each 20\,min.

\textbf{Phase tagging.} Online cycle phase $\psi(t)$ from dominant global cycle (Hilbert transform of cycle-projected component); align stimulus onsets with \emph{preceding} phase (e.g., $-500$\,ms window).

\textbf{Outcomes.} RT and accuracy as a function of phase deciles; mixed-effects models $\log RT{\sim}\beta_0{+}\beta_1\psi{+}\text{(1|subject)}$; circular-linear regressions; cluster-based permutation for phase-bins.

\textbf{Predictions.} Opposite-sign RT modulation between high-power visual/attentional vs low-power sensorimotor cycle segments; strongest effects at 300--800\,ms lags.

\textbf{Controls.} Arousal proxies (pupil, alpha), microsaccades, trial history; phase-shuffled controls.

\subsection{Study 3: Sleep Polysomnography (PSG) --- REM Desynchronization}
\textbf{Participants.} $N{=}30$ healthy sleepers (20--40\,y).

\textbf{Acquisition.} Overnight high-density EEG (128-ch), EOG, EMG, ECG; respiratory bands; sleep staged per AASM by two blinded raters. Optional MEG night ($N{=}15$) with comfortable supine protocol.

\textbf{Metrics.} Inter-network phase-locking value (PLV) and circular variance across canonical networks; cross-chain synchrony proxy from cross-frequency coupling (theta{\textrightarrow}gamma) and network HMM states; dream reports on forced awakenings.

\textbf{Predictions.} REM: preserved global cortical cycle strength but \emph{reduced cross-layer synchrony} (higher phase jitter vs NREM/wake); bizarreness ratings correlate with jitter.

\textbf{Controls.} REM latency, time of night, prior sleep debt, medication/caffeine.

\subsection{Study 4: Clinical Cohorts --- Cycle Signatures of Disorder}
\textbf{Participants.} Schizophrenia (SZ) $N{=}35$; Parkinson’s (PD) $N{=}35$; Major Depression (MDD) $N{=}35$; matched controls $N{=}45$.

\textbf{Acquisition.} Rest and task MEG as in Study 1; clinical scales (PANSS, UPDRS, HDRS); PD tested ON/OFF meds (counterbalanced).

\textbf{Hypotheses.} SZ: reduced cycle strength, elevated phase noise; PD: increased rigidity (higher coherence, lower flexibility), abnormal response to dopaminergic state; MDD: reduced repertoire (lower state entropy) and slower rates.

\textbf{Analyses.} Group mixed models on cycle metrics; mediation by clinical severity; PD ON/OFF within-subject contrasts.

\textbf{Outcomes.} Diagnostic ROC from cycle features; effect sizes for biomarker potential.

\subsection{Study 5: Causal Perturbation --- Phase-Locked Stimulation}
\textbf{Participants.} $N{=}32$ healthy adults.

\textbf{Intervention.} tACS or rhythmic TMS at individualized cycle frequency; stimulation at phase $\varphi{=}0^\circ$ vs $180^\circ$ relative to current cycle; 20\,min blocks; sham control.

\textbf{Outcomes.} Shift in cycle phase and rate; changes in RT/accuracy coupling; safety/AE monitoring.

\textbf{Predictions.} In-phase stimulation increases cycle strength/rate and behavioral phase locking; anti-phase reduces them; aftereffects scale with stimulation dose.

\textbf{Controls.} Counterbalanced order, scalp sensation ratings, blind assessment.

\subsection{Study 6: Cross-Species and Multimodal Validation}
\textbf{Rodent.} Hippocampal LFP (theta/gamma) and motor CPG recordings during memory tasks; compute cycle analogs and cross-layer synchrony; pharmacological modulation (e.g., dopaminergic agents).

\textbf{fMRI.} Human resting-state fMRI (7\,T) to validate network assignments underlying MEG cycles via co-fluctuation events; compare state-sequence asymmetries across modalities.

\subsection{Measurement Model and Derived Variables}
\textbf{Cycle strength.} Norm of antisymmetric component of lagged transition matrices (FO-asymmetry).  
\textbf{Cycle rate.} Peak frequency of global cycle component (Hilbert/FFT).  
\textbf{Phase jitter.} Circular variance of inter-network phase offsets.  
\textbf{Cross-layer synchrony.} Envelope ICMI/coherence between cortical cycle and CPG proxies (motor/phonological loops).  
\textbf{Behavioral coupling.} Slope of RT vs phase; accuracy modulation index.

\subsection{Statistical Plan and Corrections}
Frequentist: mixed-effects, cluster permutation, circular stats (Watson--Williams, Kuiper). Bayesian: hierarchical models for individual rates/strengths; ACE twin models if twin subsample available. Multiple-comparison control via FDR ($q{<}.05$). Robustness via preregistered alternative pipelines (HMM state counts, parcellations).

\subsection{Confounds and Mitigations}
Head motion, vigilance, ocular/muscle artifacts, hardware drift, diurnal effects, medication, caffeine, prior sleep. Include covariates; sensitivity analyses; replication cohort ($N{=}40$) for Studies 1--2.

\subsection{Open Science and Reproducibility}
Preprint + preregistration (OSF); code (analysis and simulations) and de-identified data released under permissive license; containerized pipelines (Docker/Singularity); method checklists per MEG‐Best‐Practices.

\subsection{Falsifiable Criteria}
RSVP--CPG is disconfirmed if: (i) FO-asymmetry $\not>{}0$ across lags after controls; (ii) behavioral performance shows no reliable phase dependence; (iii) REM shows no cross-layer desynchronization relative to wake/NREM; or (iv) phase-locked stimulation fails to modulate cycle metrics in the predicted direction.


\section{Interdisciplinary Connections: Cognitive Science Integration}
\subsection{Language and Gait}

Language is fundamentally rhythmic. Every conversation is shaped by prosody, intonation, and timing cues that resemble steps in a dance. Psycholinguistic research has shown that speakers unconsciously synchronize speech rhythms with their interlocutors, a phenomenon known as entrainment \citep{breska2017neural,levinson2016turn}. This rhythmic coordination is not decorative but constitutive: it enables smooth turn-taking, shared attention, and the rapid exchange of complex symbolic structures.

From the perspective of the RSVP--CPG framework, prosody is best understood as a linguistic gait. Just as physical walking depends on chained central pattern generators (CPGs) coordinating leg movements, conversational flow depends on chained oscillatory patterns coordinating syllables, stresses, and pauses. Each utterance is a stride; each pause is a moment of weight transfer; each prosodic contour signals readiness for the next step. Rhythm allows cognition to move forward, carrying memory and meaning along with it.

This analogy extends to syntax. Cross-linguistic studies reveal that sentence structures often align with rhythmic parsing windows, such as the $\sim$500 ms cycle of phrasal units \citep{ding2016cortical}. In RSVP--CPG terms, syntax emerges from the entrainment of mnemonic loops to cortical cycles, producing a phonological “gait” that constrains how linguistic items are sequenced. Language, in this view, is not simply a code transmitted from speaker to hearer, but a co-gaiting of oscillators across brains.

The theory also makes developmental sense. Before mastering speech, infants engage in rhythmic babbling, clapping, and gestural games like peek-a-boo. These embodied rhythms scaffold later phonological and syntactic learning. Gesture thus precedes and grounds auditory language: the first syntax is bodily, a chaining of motions into communicative sequences \citep{goldin2003language}. RSVP--CPG situates this within a broader architecture of gait: spoken language is a secondary layer, piggybacking on motor rhythms that evolved earlier for movement and coordination.

Finally, this perspective has predictive consequences. If conversational fluency depends on rhythmic entrainment, then cortical cycle phase should predict the ease of turn-taking and the alignment of prosodic contours. Disorders of speech and language, from stuttering to aphasia, may reflect disruptions in cycle synchrony across CPG chains. In practical applications, speech therapy and language learning might be enhanced by rhythmic entrainment protocols—treating language not as code to be decoded, but as gait to be learned and stabilized.

Language, then, is not an exception to the RSVP--CPG framework but one of its clearest exemplars. The gait of thought is also the gait of speech: a perpetual oscillatory progression that binds memory, meaning, and social coordination into rhythmic motion.


\subsection{Social Cognition}

Interpersonal synchrony facilitates empathy, with cultural variations in entrainment patterns.

\subsection{Learning and Memory}

New gaits develop through practice, with consolidation stabilizing rhythmic patterns.

\section{Broader Scientific Connections}

\subsection{Evolutionary Biology}

Central pattern generators (CPGs) predate the cortex by hundreds of millions of years, controlling locomotion, respiration, and other rhythmic functions across vertebrates and invertebrates. Their conservation across species suggests that rhythmic oscillation is a core biological strategy for coordinating distributed systems \citep{grillner2006biological}. From an evolutionary perspective, RSVP--CPG argues that cortical cycles did not emerge de novo but represent a cortical reuse of these ancient rhythmic mechanisms.

Comparative evidence supports this view. Invertebrates such as crustaceans exhibit highly structured CPG circuits for locomotion and feeding. In mammals, spinal CPGs generate walking patterns even in the absence of cortical control. The cortical cycles identified by Van Es et al. \citep{van2025cortical} may thus be understood as higher-order descendants of these ancestral oscillators: rather than driving limbs directly, they sequence perceptual, mnemonic, and attentional “steps.”

Ontogenetic development mirrors this phylogenetic layering. Infants exhibit rhythmic behaviors—suckling, babbling, rocking—long before voluntary language or abstract reasoning. These early cycles scaffold later cognitive gaits, just as locomotor CPGs scaffold the emergence of complex motor repertoires. RSVP--CPG interprets this continuity as evidence that cortical oscillations recruit and extend ancient motor architectures for cognitive purposes.

This evolutionary account has testable implications. If cortical cycles are reused motor gaits, then (i) genetic variation in motor CPGs should covary with cycle rate heritability, (ii) cross-species comparisons should reveal homologous cycle periods between locomotor rhythms and cortical rhythms, and (iii) evolutionary transitions in communication (e.g., birdsong, primate gestural repertoires) should align with diversification of CPG-like cortical dynamics. Such findings would ground RSVP--CPG not only in contemporary neuroscience but also in a deep evolutionary trajectory.

In this light, the gait metaphor is not incidental. It reflects an evolutionary throughline: from ancestral locomotor oscillations, to infant rhythmic behaviors, to adult cortical cycles orchestrating cognition. Consciousness as gait is thus both a phylogenetic inheritance and a cortical innovation.

\section{Category-Theoretic Interpretation of RSVP--CPG}

The RSVP--CPG framework can be reformulated in categorical terms to highlight its structural coherence.  
We treat cortical cycles, mnemonic proxies, and CPG chains as \emph{objects} in a category $\mathcal{C}$, with morphisms representing lawful transitions between rhythmic states.  

\subsection{Objects and Morphisms}
\begin{itemize}
  \item \textbf{Objects}: Cycle attractors, memory loops, and gait modules.
  \item \textbf{Morphisms}: Directed asymmetric transitions, modeled as entropy-smoothing maps $\lambda : X \to Y$ between oscillatory states.
\end{itemize}

The cortical cycle layer then becomes a diagram in $\mathcal{C}$, with arrows encoding transition probabilities. Limit cycles correspond to categorical \emph{endofunctors} $F : \mathcal{C} \to \mathcal{C}$ whose iterates preserve structure.  
Lyapunov stability is expressed as the existence of natural transformations $\eta: 1_{\mathcal{C}} \Rightarrow F$ that witness coherence across iterations.

\subsection{Compositional Gait}
Cognition as gait can thus be understood as the composition of morphisms along oscillator chains. A memory proxy $\Pi_m$ is not a stored entity but a recurrent morphism in the category of gait transitions. Dream desynchronization corresponds to a failure of naturality: functorial action ceases to commute across chains, producing “paradoxical” compositions.

\section{Sheaf-Theoretic Interpretation of RSVP--CPG}

Where category theory captures compositionality, sheaf theory formalizes the local-to-global structure of cyclical cognition.

\subsection{Open Covers of Cognitive Space}
Let $X$ be the cortical manifold of functional networks. A cover $\{U_i\}$ represents local cycle phases (visual, mnemonic, motor, default-mode).  
Sections $s_i \in \mathcal{F}(U_i)$ correspond to local rhythmic solutions—stable attractors of RSVP dynamics within that network.

\subsection{Gluing Conditions}
Sheaf gluing captures synchronization:
\begin{itemize}
  \item \textbf{Consistency}: On overlaps $U_i \cap U_j$, local sections must agree up to phase shifts.
  \item \textbf{Gluing}: If local cycle phases match on overlaps, there exists a global gait $s \in \mathcal{F}(X)$.
\end{itemize}

In waking cognition, gluing succeeds, yielding coherent global gait. In REM, precision weighting collapses, and gluing fails: local sections persist but no consistent global section exists, producing dreamlike patchworks.

\subsection{Markov Blankets as Sheaves}
Markov blanket partitions correspond to sub-sheaves: self, other, and world rhythms are local sections glued differently depending on entrainment. Skeuomorphic affordances are extendable sheaf morphisms: design succeeds when tool dynamics extend an existing section of the body’s own sheaf of gaits.

\section{Homotopy- and Type-Theoretic Perspectives on RSVP--CPG}

The categorical and sheaf-theoretic formalisms of RSVP--CPG can be deepened by appeal to homotopy type theory (HoTT) and higher categories, which allow rhythmic cognition to be described as a space of equivalence classes and higher-order invariants.

\subsection{Cycles as Homotopy Classes}
Each cortical cycle can be represented as a closed path $\gamma: S^1 \to X$ in the cortical state space $X$.  
Two cycles $\gamma_1, \gamma_2$ are equivalent if there exists a homotopy $H: S^1 \times [0,1] \to X$ deforming one into the other without leaving the attractor basin.  
Thus, cycles belong to $\pi_1(X)$, the fundamental group of cortical rhythm space, capturing the idea that “gait” is a topological invariant of neural dynamics.

\subsection{Higher-Categorical Synchrony}
CPG chains form higher-categorical structures:  
\begin{itemize}
  \item \textbf{1-morphisms}: transitions between local oscillatory states, as in the categorical view.  
  \item \textbf{2-morphisms}: synchronizations between entire chains, i.e. transformations of transformations.  
\end{itemize}
REM desynchronization corresponds to the breakdown of coherence at the level of 2-morphisms: local morphisms persist but higher coherence laws fail.

\subsection{Type-Theoretic Encoding}
In a type-theoretic framework:
\begin{itemize}
  \item A cycle is a \emph{path type}, witnessing the equality of states across time.  
  \item A memory proxy is a dependent type, whose validity depends on the maintenance of a phase relation.  
  \item Dream recombination arises when type-checking fails: paths no longer commute, producing non-canonical constructions that are nevertheless locally valid.  
\end{itemize}

\subsection{Persistence and Topological Data Analysis}
Using persistent homology, RSVP--CPG predicts that cortical rhythms form stable barcodes: intervals of homology classes that persist across scales. Waking cognition corresponds to long-persistence features (stable gaits), while REM corresponds to short-lived or noisy homological bars, reflecting unstable coordination.

\subsection{Summary}
By embedding RSVP--CPG into homotopy type theory and persistent topology, we formalize gait as a higher-order invariant: cognition is not only rhythmic motion but a structure preserved up to deformation.  
This connects phenomenology (dreams, agency, memory loops) with mathematics of paths, equivalences, and higher symmetries, offering a rigorous foundation for the claim that consciousness is the gait of thought.

\subsection{Empirical Predictions from the Homotopy View}

The higher-categorical formalization of RSVP--CPG yields specific, testable predictions:

\paragraph{1. Persistence barcodes in MEG.}  
If cortical cycles are topological invariants, then topological data analysis (TDA) applied to MEG-derived state-space trajectories should reveal persistent 1-dimensional homology classes (loops) corresponding to cortical gaits.  
\emph{Prediction:} Waking states yield long-persistence loops; REM yields shorter-lived, noisy barcodes reflecting desynchronization.

\paragraph{2. Homotopy failures in REM.}  
Dream bizarreness is modeled as failure of path compositions: cycles no longer commute in the fundamental groupoid of cortical state space.  
\emph{Prediction:} REM MEG/fMRI trajectories show higher rates of non-commutative path transitions, measurable as increased entropy in cycle ordering.

\paragraph{3. Cognitive load as homotopy contraction.}  
Under high working-memory load, mnemonic proxy loops are forced into tighter alignment with cortical cycles, reducing homotopic diversity.  
\emph{Prediction:} Cycle diversity (number of distinguishable loops in $\pi_1$ estimates) decreases as working-memory demands increase.

\paragraph{4. Interpersonal synchrony as higher morphisms.}  
Dyadic interactions (conversation, joint action) can be modeled as 2-morphisms aligning two individuals’ cortical cycles.  
\emph{Prediction:} Hyperscanning EEG/MEG reveals higher coherence in phase-path alignments during successful joint tasks than during asynchronous interaction.

\paragraph{5. Tool learning as persistent homology.}  
Skeuomorphic interfaces entrain existing gaits, yielding long-lived topological features in user–tool interaction space.  
\emph{Prediction:} Motion-tracking data during tool learning shows faster convergence to persistent homology classes when skeuomorphic affordances are present.

\paragraph{Summary.}  
These predictions connect abstract invariants (homotopy, homology, higher morphisms) with measurable neurodynamic signatures. RSVP--CPG thus links categorical semantics with experimental protocols, offering a rare bridge between higher mathematics and empirical neuroscience.

\subsection{Methods Sketch: Topological and Groupoid Analyses of Cortical Cycles}

\paragraph{Data and Preprocessing.}
MEG preprocessing follows Sec.~\ref{sec:empirical}: SSS/MaxFilter, ICA/SSP artifact removal, 1--100\,Hz bandpass, source projection (LCMV) to Schaefer-200 ROIs. For fMRI, use HCP-style preprocessing and nuisance regression; extract ROI time series (TR-matched).

\paragraph{State Trajectories.}
Construct low-dimensional trajectories $X(t)\in\mathbb{R}^d$ via:
(i) HMM state-occupancy vectors (dim.\,$d{=}$\#states); or
(ii) Takens delay-embedding on principal components of ROI power:
\[
X(t) = \big[ y(t), y(t-\tau), \dots, y(t-(m-1)\tau) \big],
\]
with $y(t)$ the first 3--5 PCs of band-limited envelopes; choose $\tau$ by first minimum of average mutual information; $m$ via false-nearest neighbors (target $d{\le}8$).

\paragraph{Distance Geometry and Denoising.}
Compute pairwise distances $D_{ij}=\|X(t_i)-X(t_j)\|_2$. Optionally, apply diffusion maps or local tangent space alignment to stabilize geometry; sparsify by $k$-NN (e.g., $k{=}20$) to reduce noise.

\paragraph{Persistent Homology (PH).}
Build a Vietoris--Rips filtration on $(X,D)$; compute persistence diagrams $\mathsf{Dgm}_k$ (Ripser/GUDHI) for $k=0,1,2$; summarize by total persistence $\mathsf{TP}_k$, max bar length $\mathsf{MBL}_k$, count of long bars $N_k(\epsilon)$, and bottleneck/Wasserstein distances between conditions (wake, NREM, REM). Primary target is $k{=}1$ loops (gait cycles).

\emph{Predictions:} Wake $\Rightarrow$ larger $\mathsf{MBL}_1$ and $N_1(\epsilon)$; REM $\Rightarrow$ shorter, noisier $H_1$ bars; NREM $\Rightarrow$ reduced $\mathsf{TP}_1$, increased $H_0$ persistence (fragmentation).

\paragraph{Sliding-Window PH.}
Use windows of 30--60\,s with 50\% overlap; track $\mathsf{MBL}_1(t)$ and $N_1(t)$; correlate with vigilance proxies (alpha power, pupil) and behavioral blocks. Permutation tests on time-aligned surrogates (phase-randomized).

\paragraph{Cycle Extraction and Groupoid Reconstruction.}
From HMM transitions, estimate lagged transition matrices $T^{(\Delta)}$; extract cyclic orderings by antisymmetric component $A^{(\Delta)}=(T^{(\Delta)}-(T^{(\Delta)})^\top)/2$. Build a directed multigraph $\mathcal{G}$ of recurrent cycles. Define partial compositions of paths to form a \emph{fundamental groupoid} $\Pi(\mathcal{G})$: objects are recurrent states; morphisms are homotopy classes of admissible paths (admissibility via energy/entropy thresholds).

\emph{Non-commutativity metric.} For pairs of frequently observed loops $\alpha,\beta$, compute
\[
\Delta_{\mathrm{nc}}=\mathbb{E}\big[ d\big(\alpha\!\circ\!\beta,\ \beta\!\circ\!\alpha\big)\big],
\]
where $d$ is edit or dynamic time-warp distance between state sequences projected back to $X$. \emph{Prediction:} $\Delta_{\mathrm{nc}}$ increases in REM relative to wake.

\paragraph{Homology Stability and Barcodes as Biomarkers.}
Use bootstrap subsampling to assess barcode stability (bottleneck variance). Define an index
\[
\mathcal{B}_1=\mathrm{MBL}_1 \times N_1(\epsilon)
\]
as a compact gait persistence score. \emph{Prediction:} $\mathcal{B}_1$ decreases with SZ symptom severity; increases with dopaminergic ON state in PD if rigidity \emph{decreases} (or decreases if rigidity \emph{increases}); correlates with RSVP cycle strength.

\paragraph{Hyperscanning and Higher Morphisms.}
Collect dual-EEG/MEG during joint tasks (synchronous tapping, conversational turn-taking). Build individual trajectories $X^{(1)}(t), X^{(2)}(t)$ and a joint trajectory $X^{(12)}(t){=}[X^{(1)}(t),X^{(2)}(t)]$. Compute PH on joint space and compare against the disjoint union baseline (shuffle partners). Define a 2-morphism coherence index as the excess $H_1$ persistence in joint vs shuffled pairs. \emph{Prediction:} successful coordination $\Rightarrow$ higher joint persistence and lower $\Delta_{\mathrm{nc}}$ across partners.

\paragraph{Software and Parameters.}
PH: Ripser, GUDHI, giotto-tda; embeddings: scikit-learn, PyEMD; groupoid ops: custom Python (NetworkX) with unit tests. Defaults: $k$-NN=20, window=45\,s, step=22.5\,s, $\epsilon$ at 75th percentile of $H_1$ bar lengths. All parameter choices preregistered with sensitivity analyses.

\paragraph{Statistics.}
Within-subject contrasts (wake vs REM vs NREM): linear mixed models on $\mathsf{MBL}_1$, $\mathcal{B}_1$, $\Delta_{\mathrm{nc}}$; multiple-comparison control via FDR; circular-linear models for phase covariates. Surrogate testing (phase-randomized, time-scrambled) for null distributions of barcodes.

\paragraph{Quality Control.}
Artifact sensitivity checks (repeat PH after excluding high-motion epochs); replicate across embeddings (HMM vs Takens) and distances ($\ell_2$ vs cosine). Public release of code/notebooks and containerized environments for reproducibility.


\section{Conclusion: Toward a Rhythmic Science of Consciousness}

The evidence reviewed here suggests that large-scale cortical networks do not activate randomly or in purely stochastic bursts, but in robust, cyclical sequences with stable ordering, asymmetric transitions, and heritable rates \citep{van2025cortical}. These findings resonate with long-standing philosophical intuitions—from Aristotle’s peripatetic pedagogy to Merleau-Ponty’s phenomenology—that thinking is inseparable from motion. They also converge with embodied cognition research in music \citep{cox2016music} and spatial reasoning \citep{tversky2019mind}, both of which emphasize rhythm and movement as the deep structure of thought.

By formalizing these insights, the RSVP--CPG framework advances a unified claim: consciousness is gait. Cortical cycles operate as chained central pattern generators, memory persists as rhythmic proxy loops entrained to those cycles, and paradoxical sleep reflects their desynchronization. This reconceptualization reframes cognition not as equilibrium (Active Inference) nor as spotlight (Global Workspace Theory), but as rhythmic descent through entropic landscapes—perpetually falling forward into the next step.

The framework integrates across levels:

\begin{itemize}[noitemsep,topsep=0pt]
\item \textbf{Neurocomputational.} Mathematical models of coupled oscillators demonstrate how asymmetric transitions, Lyapunov-stable limit cycles, and entropy pacing produce the observed cortical cycles. Bayesian analyses recover cycle order, rate, and heritability from empirical data, validating the generative model.
\item \textbf{Cognitive and Clinical.} The gait metaphor explains working memory as proxy loops, theory of mind as entrainment to another’s movements, and REM bizarreness as cross-chain desynchronization. It generates testable biomarkers, from persistence scores in persistent homology to non-commutativity metrics in cycle groupoids.
\item \textbf{Philosophical.} RSVP--CPG deepens the lineage from peripatetic walking to modern enactivism. It aligns with process philosophy (Whitehead, Bergson), phenomenology (Husserl, Heidegger), and pan-embodied views of cognition, offering rhythm as the missing category linking self, other, and world.
\item \textbf{Formal.} Category-theoretic and sheaf-theoretic interpretations reframe cycles as coverings, groupoids, and homotopy classes. Cognitive operations become gluing problems, and REM desynchronization emerges as failure of colimits in sheaf assignments. Topological data analysis provides practical measures of these structures in MEG data.
\item \textbf{Applied.} Skeuomorphic affordances in tool design, AI architectures with oscillatory priors, and clinical interventions targeting cycle synchrony exemplify the translational potential of the framework.
\end{itemize}

What emerges is not simply another metaphor for consciousness, but a candidate substrate: an oscillatory skeleton that explains why phenomenology is rhythmic, why cognition remains temporally ordered, and why disruption of cycles produces both psychiatric symptoms and dreamlike distortions. By locating cognition in non-equilibrium limit cycles, RSVP--CPG bridges the gap between physical dynamics and lived experience.

Future work must extend this account empirically, through multimodal neuroimaging and hyperscanning, and formally, through deeper integration with Active Inference, IIT, and dynamic core models. Yet the direction is clear. To think is to move, and to move is to think. The rhythms of cortical cycles are not incidental correlates, but the gait of consciousness itself: a perpetual dance of self, other, and world, paced by entropy and sustained by motion.

\subsection{Complex Systems Science}

The synchronization of neural oscillations parallels the general principles of complex systems, where global order emerges from local interactions. In physics, phenomena such as coupled oscillators, spin glasses, and flocking dynamics show how simple local rules generate macroscopic coherence. The RSVP--CPG framework interprets cortical cycles through this lens: asymmetric transitions and recursive pacing allow the brain to self-organize into stable limit cycles without central coordination. This situates cortical dynamics within a broader family of critical phenomena, poised between order and disorder. Importantly, this perspective suggests that cognition inherits the robustness and adaptability characteristic of other complex systems, explaining how the brain resists perturbations yet remains flexible to new information. The implication is that consciousness may not be an exceptional process, but rather an instance of universal self-organizing principles.

\subsection{Applied Fields}

Rhythmic entrainment offers promising applications across educational, therapeutic, and technological domains. In education, structuring lessons with rhythmic pacing—whether through music, movement, or temporally scaffolded activities—enhances retention and recall by aligning cognitive gaits with external rhythms. In clinical contexts, disorders such as dyslexia, ADHD, and Parkinson’s disease can be reframed as disruptions of cognitive gait, opening avenues for rhythm-based interventions including metronomic training or neurostimulation. In human–computer interaction, skeuomorphic affordances leverage entrainability, allowing interfaces to “click” with users’ motor expectations. Extending RSVP--CPG into AI design suggests architectures where oscillatory priors structure inference, improving robustness and preventing catastrophic forgetting. These applications underscore that the rhythmic substrate of cognition is not an abstraction but a practical lever for design and therapy.

\section{Technical Expansions: Advanced Mathematical Tools}

\subsection{Information Geometry}

Information geometry provides a powerful lens for formalizing the structure of cycle space. Each cortical state can be modeled as a probability distribution over network configurations, with cycles tracing geodesics on the statistical manifold. Fisher information metrics quantify the sensitivity of cycles to perturbations, while divergence measures capture the “distance” between distinct gaits of cognition. RSVP--CPG thus frames cognition as motion across an information manifold, where stability corresponds to geodesic curvature and learning corresponds to re-shaping the underlying metric. This geometric framing connects directly to Active Inference and variational principles, but emphasizes oscillatory trajectories rather than equilibrium points, offering a richer vocabulary for non-stationary processes such as REM desynchronization.

\subsection{Data Analysis Tools}

To test RSVP--CPG empirically, specialized algorithms for detecting and characterizing cycles are essential. Cycle detection can be achieved through phase-space reconstruction and recurrence quantification analysis, while phase coupling is measured via metrics such as phase-locking value (PLV), cross-frequency coupling, and topological persistence of loops in time-series data. Recent advances in topological data analysis (TDA) provide additional tools: persistent homology detects stable loops across scales, while Mapper algorithms reveal transitions between states. Integrating these approaches enables researchers to identify whether empirical neural data conforms to the Lyapunov-stable cycles predicted by RSVP--CPG. Open-source pipelines combining MEG preprocessing, phase estimation, and persistence barcodes will make such analyses reproducible and widely accessible.

\section{Potential Challenges and Solutions: Common Pitfalls}

One challenge for RSVP--CPG is the risk of over-interpretation: oscillatory patterns are ubiquitous in the brain, but not all oscillations are signatures of cognition or consciousness. To avoid conflating correlation with causation, claims must remain anchored in testable predictions and falsifiable hypotheses. Another pitfall is uncritical reliance on metaphor. While the gait analogy is fruitful, it must be continuously translated into rigorous mathematics and empirical measures. A further challenge is disciplinary isolation: to gain traction, RSVP--CPG must engage not only neuroscientists but also philosophers, physicists, and AI researchers. Solutions include: (i) iterative model refinement through open datasets, (ii) explicit cross-validation with competing theories such as GWT and AIF, and (iii) clear specification of limitations—acknowledging phenomena the framework cannot yet explain. By combining ambition with methodological caution, RSVP--CPG can advance without overstating its reach.


\section{Conclusion: Oscillatory Attractors and the Rhythmic Science of Consciousness}

Recent evidence that large-scale cortical networks exhibit robust cyclical activation patterns \citep{van2025cortical} underscores the centrality of oscillatory attractors in the organization of cognition. These cycles are not random fluctuations but structured, asymmetric sequences with stable ordering and heritable rates. Such findings reveal an underlying skeleton of cortical dynamics that existing theories only partially explain. Global Workspace Theory provides a phenomenological framework for conscious access but does not specify why transitions are rhythmic. Active Inference offers a mathematically precise account of predictive regulation but presupposes equilibrium-seeking dynamics, leaving paradoxical states such as REM sleep underexplained.

The RSVP--CPG framework situates cortical cycles as the fundamental substrate: cognition emerges from chained oscillatory processes that enforce ordered transitions across functional networks. These cycles operate as non-equilibrium limit cycles, maintaining progression through recursive smoothing of entropic constraints. Our Bayesian reformulation of the Van Es et al. findings formalizes these dynamics by estimating global cycle order, asymmetry, and rate as parameters of coupled oscillator systems. This formulation generates empirically testable predictions: cycle strength will covary with working memory precision, cycle desynchronization will track REM-specific dream phenomenology, and heritability of cycle rate will reflect genetic modulation of oscillator coupling.

But RSVP--CPG extends beyond neurocomputation. The framework integrates multiple levels of analysis:

\begin{itemize}[noitemsep,topsep=0pt]
\item \textbf{Neurocomputational.} Mathematical models of coupled oscillators demonstrate how asymmetric transitions, Lyapunov-stable limit cycles, and entropy pacing reproduce the observed cortical cycles. Bayesian analyses recover cycle order, rate, and heritability, validating RSVP--CPG as a generative model of MEG data.
\item \textbf{Cognitive and Clinical.} The gait metaphor explains working memory as rhythmic proxy loops, theory of mind as entrainment to others’ movements, and REM bizarreness as cross-chain desynchronization. These insights yield testable biomarkers and therapeutic interventions, from cycle coherence measures to stimulation protocols targeting synchrony.
\item \textbf{Philosophical.} RSVP--CPG resonates with the peripatetic claim that to think is to move. It aligns with process philosophy (Whitehead, Bergson), phenomenology (Husserl, Heidegger), and embodied cognition (Cox, Tversky), offering rhythm as the missing category linking self, other, and world. Consciousness emerges not as spotlight or equilibrium but as gait: a perpetual fall forward through time.
\item \textbf{Formal.} Category-theoretic and sheaf-theoretic interpretations reframe cortical cycles as coverings, groupoids, and homotopy classes. Cognitive operations become gluing problems; REM desynchronization corresponds to failed colimits. Topological data analysis offers practical tools to measure these structures in empirical data.
\item \textbf{Applied.} Skeuomorphic affordances in interface design, oscillatory priors in AI architectures, and clinical interventions for cycle regulation illustrate the translational reach of RSVP--CPG. From human–computer interaction to psychiatry, the rhythmic substrate provides a new axis for design and diagnosis.
\end{itemize}

What emerges is not merely another metaphor for consciousness, but a candidate substrate: an oscillatory skeleton that explains why phenomenology is rhythmic, why cognition remains temporally ordered, and why disruption of cycles produces psychiatric symptoms and dreamlike distortions. By identifying cortical cycles as Lyapunov-stable attractors rather than incidental correlates, RSVP--CPG provides the missing oscillatory substrate that unifies conscious access, predictive regulation, and dream phenomenology.

Future work must extend this account empirically, through multimodal neuroimaging and hyperscanning, and formally, through deeper integration with Active Inference, Integrated Information Theory, and dynamic core models. Yet the direction is clear. To think is to move, and to move is to think. The rhythms of cortical cycles are not incidental correlates but the gait of consciousness itself: a perpetual dance of self, other, and world, paced by entropy and sustained by motion.

\newpage
\appendix
\section{Mathematical Formalization and Stability Analysis}
\label{app:math}

\subsection{Field Equations of RSVP--CPG}

We model cortical dynamics as coupled scalar, vector, and entropy fields defined on a spatial domain 
$\Omega \subset \mathbb{R}^3$ with time parameter $t$:
\begin{align}
\partial_t \Phi(x,t) &= -\alpha \Phi(x,t) + \nabla \cdot \mathbf{v}(x,t) - \lambda S(x,t) + \eta_\Phi(x,t), \label{eq:phi}\\
\partial_t \mathbf{v}(x,t) &= -\beta \mathbf{v}(x,t) - \nabla U(\Phi(x,t)) + \gamma \nabla S(x,t) + \eta_v(x,t), \label{eq:v}\\
\partial_t S(x,t) &= \delta \nabla \cdot \mathbf{v}(x,t) - \kappa S(x,t) + f(\Phi,\mathbf{v}) + \eta_S(x,t). \label{eq:S}
\end{align}

Here:
\begin{itemize}
\item $\Phi$ is a scalar density field (neural activation potential),
\item $\mathbf{v}$ is a vector flow field (cortical CPG propagation),
\item $S$ is an entropy field (local cycle disorder),
\item $U(\Phi)$ is an effective potential (e.g.\ sigmoid or quartic),
\item $\alpha,\beta,\lambda,\gamma,\delta,\kappa >0$ are coupling constants,
\item $\eta_\Phi, \eta_v, \eta_S$ represent stochastic perturbations.
\end{itemize}

Equations \eqref{eq:phi}--\eqref{eq:S} define the \textbf{RSVP--CPG field system}. 
They capture the intuition that scalar density feeds vector flow, vector flow modulates entropy, and entropy 
feeds back into scalar suppression.

\subsection{Limit Cycles and CPG Dynamics}

To illustrate gait-like oscillations, consider a spatially homogeneous reduction:
\begin{align}
\dot{\phi} &= -\alpha \phi + \mathbf{v}, \\
\dot{v} &= -\beta v - u(\phi) + \gamma s, \\
\dot{s} &= \delta v - \kappa s + g(\phi,v),
\end{align}
where $(\phi,v,s)$ are scalar reductions of $(\Phi,\mathbf{v},S)$ and $u, g$ are nonlinear coupling functions.

This three-dimensional ODE system supports limit cycles under Hopf bifurcation conditions. 
Specifically, for $u(\phi) = \mu \phi$ and $g(\phi,v) = \nu \phi v$, 
one finds oscillatory solutions when $\delta \gamma > \alpha \beta$.

\subsection{Lyapunov Stability}

Define a candidate Lyapunov function:
\begin{equation}
V(\phi,v,s) = \frac{1}{2}\left( \phi^2 + v^2 + s^2 \right).
\end{equation}
Then,
\begin{align}
\dot{V} &= \phi \dot{\phi} + v \dot{v} + s \dot{s} \\
&= -\alpha \phi^2 - \beta v^2 - \kappa s^2 + (\text{cross-terms from coupling}). 
\end{align}

If parameters satisfy
\[
\alpha, \beta, \kappa > 0 \quad \text{and} \quad \delta \gamma \approx \alpha \beta,
\]
the system admits bounded oscillations: trajectories are Lyapunov stable 
around a limit cycle rather than converging to a fixed point. 
This formalizes the intuition that cortical dynamics are perpetually non-equilibrium but stable.

\subsection{Entropy Functional and Variational Principle}

Define a cycle entropy functional over network states:
\begin{equation}
\mathcal{S}[p] = -\int_\Omega p(x) \log p(x)\, dx,
\end{equation}
where $p(x)$ is the probability distribution of network states. 

The RSVP--CPG dynamics can be interpreted as gradient flows minimizing a 
\textbf{free-energy-like functional}:
\begin{equation}
\mathcal{F}[p] = D_{\mathrm{KL}}(p \,\|\, q) + \Lambda \cdot \mathcal{S}[p],
\end{equation}
where $q$ is a reference (equilibrium) distribution and $\Lambda$ sets the 
entropy-smoothing tradeoff. 

Thus the RSVP--CPG field equations can be derived as flows that simultaneously 
\textbf{stabilize oscillations (Lyapunov cycles)} and \textbf{minimize entropy gradients}, 
consistent with Active Inference formulations.

\section{Experimental Protocols and Predictions}
\label{app:empirical}

This appendix outlines concrete empirical tests of the RSVP--CPG framework, 
designed to validate cyclical cortical dynamics and their role in cognition.  

\subsection{Core Predictions}

\begin{enumerate}
    \item \textbf{Cortical cycles exhibit stable limit cycles.}  
    Large-scale MEG/EEG networks should reveal asymmetric transitions 
    forming recurrent cycles, as reported in \citet{vanes2025}.  

    \item \textbf{Memory loops entrain to cortical gait.}  
    Behavioral measures of working memory should covary with cross-frequency coupling 
    between theta and gamma oscillations in MEG.  

    \item \textbf{REM sleep reflects desynchronized CPG chains.}  
    Polysomnography should reveal increased phase jitter across cortical networks 
    during paradoxical sleep compared to NREM.  

    \item \textbf{Clinical disorders reflect gait disruptions.}  
    - Schizophrenia: fragmented cycles and reduced cycle coherence.  
    - Parkinson’s: excessive rigidity in CPG coupling.  
    - Depression: narrowed repertoire of cyclical transitions.  
\end{enumerate}

\subsection{Study 1: MEG Detection of Cortical Cycles}

\textbf{Participants:} $N=40$ healthy adults (balanced gender, ages 18--35).  
\textbf{Acquisition:} 306-channel Elekta Neuromag MEG, 1 kHz sampling, 30 min resting-state + cognitive tasks.  
\textbf{Analysis:}  
\begin{itemize}
    \item Hidden Markov Model (HMM) inference of network states.  
    \item Transition matrices tested for asymmetry and cycle decomposition.  
    \item Statistical test: permutation analysis with cluster correction, $\alpha = 0.05$.  
\end{itemize}  
\textbf{Prediction:} Cyclical network transitions with robustness across participants.  

\subsection{Study 2: Memory Loops and Working Memory}

\textbf{Participants:} $N=30$ adults, $N=30$ adolescents.  
\textbf{Task:} n-back task with visual stimuli, 20 min duration.  
\textbf{Acquisition:} Simultaneous MEG and EEG.  
\textbf{Analysis:}  
\begin{itemize}
    \item Cross-frequency coupling (theta $\times$ gamma).  
    \item Correlation of coupling strength with task accuracy.  
    \item Power analysis: sample size ensures detection of $r=0.35$ with 80\% power.  
\end{itemize}  
\textbf{Prediction:} Stronger loop entrainment correlates with improved working memory.  

\subsection{Study 3: REM Desynchronization}

\textbf{Participants:} $N=20$ healthy sleepers.  
\textbf{Acquisition:} Overnight polysomnography with high-density EEG (128-channel).  
\textbf{Analysis:}  
\begin{itemize}
    \item Phase-locking value (PLV) across networks.  
    \item Jitter analysis: variance of inter-network phase offsets.  
    \item Comparison: REM vs NREM vs wake.  
\end{itemize}  
\textbf{Prediction:} REM characterized by cycle desynchronization (higher PLV variance).  

\subsection{Study 4: Clinical Populations}

\textbf{Groups:}  
\begin{itemize}
    \item Schizophrenia patients ($N=25$),  
    \item Parkinson’s disease patients ($N=25$),  
    \item Major depressive disorder patients ($N=25$),  
    \item Controls ($N=30$).  
\end{itemize}  
\textbf{Acquisition:} Resting-state MEG + task-based EEG.  
\textbf{Predictions:}  
\begin{itemize}
    \item Schizophrenia: fragmented network cycles, reduced entropy smoothing.  
    \item Parkinson’s: exaggerated rigidity, reduced flexibility in phase transitions.  
    \item Depression: limited attractor repertoire, reduced cycle diversity.  
\end{itemize}  

\subsection{Multi-Modal Validation}

\begin{itemize}
    \item fMRI: resting-state connectivity patterns compared to cycle structure.  
    \item Single-unit animal recordings: rodent hippocampal theta/gamma cycles as proxies.  
    \item Pharmacological perturbation: dopamine modulation in Parkinson’s to test CPG rigidity.  
\end{itemize}  

\subsection{Data Sharing and Reproducibility}

All code and anonymized datasets will be made available in an open repository, 
ensuring reproducibility and cross-lab validation.  

\section{Comparative Framework}
\label{app:comparative}

This appendix situates RSVP--CPG within the broader landscape of consciousness theories. 
Rather than a tabular summary, each comparison is presented as a narrative contrast, 
emphasizing conceptual metaphors, mechanisms, and explanatory scope.  

\subsection{Global Workspace Theory (GWT)}
Global Workspace Theory casts cognition as a competition among modular processes, 
with the ``winner'' gaining access to a broadcast channel of conscious awareness.  
The metaphor is theatrical: a stage, a spotlight, an audience. RSVP--CPG differs fundamentally: 
it does not rely on ignition events or discrete access, but on rhythmic sequencing.  
Where GWT focuses on access and visibility, RSVP--CPG emphasizes gait and oscillation.  
Both theories agree that coordination across distributed networks is crucial, 
but RSVP--CPG replaces the spotlight metaphor with a locomotor one: 
progress comes from ordered cycles, not from central broadcasting.  

\subsection{Active Inference (AIF)}
Active Inference models cognition as Bayesian inference through the minimization 
of variational free energy. Its metaphor is equilibrium: a system constantly rebalancing 
its predictions against sensory evidence. RSVP--CPG adopts some of this formal structure 
but rejects equilibrium as the final aim. Cognition, in this view, is inherently 
non-equilibrium: stable but forward-driven oscillations akin to gait.  
Where AIF posits convergence, RSVP--CPG emphasizes renewal: the brain advances by 
falling forward into the next cycle, never reaching rest.  

\subsection{Integrated Information Theory (IIT)}
Integrated Information Theory explains consciousness in terms of the quantity and 
structure of integrated information (\(\Phi\)). Its metaphor is informational geometry: 
consciousness corresponds to irreducible causal structures. RSVP--CPG departs by 
locating the foundation not in static integration, but in dynamic sequencing.  
IIT’s $\Phi$ captures informational richness at a snapshot; RSVP--CPG instead focuses 
on the rhythmic passage from one state to another. From this perspective, integration 
is a byproduct of cyclical entrainment, not its essence.  

\subsection{Dynamic Core Hypothesis}
The Dynamic Core Hypothesis emphasizes rapid reconfigurations of thalamo-cortical 
loops to account for conscious states. RSVP--CPG agrees on the centrality of 
recurrent dynamics but interprets them differently. For RSVP--CPG, what matters 
is not transient reconfiguration per se, but the existence of stable cycles across 
time that structure cognitive operations as gaits. The ``core'' is thus a rhythm, 
not a fleeting coalition.  

\subsection{Attention Schema Theory (AST)}
Attention Schema Theory posits that consciousness arises when the brain constructs 
a simplified model of its own attentional processes. RSVP--CPG can complement this view: 
the construction of an attention schema may be one specific function carried by the 
cyclical gait. Yet RSVP--CPG reframes the explanatory ground: rather than a model 
of attention giving rise to consciousness, oscillatory gait provides the substrate 
on which such modeling becomes possible.  

\subsection{Predictive Processing}
Predictive Processing overlaps closely with AIF, presenting the brain as a prediction 
engine. RSVP--CPG accepts the importance of prediction but shifts the metaphor: 
prediction is not a convergence toward stasis but a coordination of rhythmic progressions. 
The ``errors'' in prediction are smoothed not by elimination but by being woven into 
the ongoing gait.  

\subsection{RSVP--CPG as Synthesis}
RSVP--CPG reframes the central question: not how information is broadcast, 
or minimized, or integrated, but how cognition advances through ordered oscillations.  
Consciousness, on this account, is not a light, not an equilibrium, not a quantity, 
but a rhythm --- a gait that structures both perception and action in time.  
Other theories capture partial aspects of this rhythm, but RSVP--CPG proposes 
that the oscillatory skeleton is the primary substrate on which broadcasting, 
prediction, integration, and modeling all depend.  

\section{Interdisciplinary Connections}
\label{app:interdisciplinary}

The RSVP--CPG framework does not stand in isolation as a neurodynamical theory.  
It intersects with multiple domains across cognitive science, evolutionary biology, 
and applied research. This appendix sketches several of these connections.  

\subsection{Language as Rhythmic Entrainment}
Spoken language is layered upon rhythmic substrates.  
Prosody, turn-taking, and gestural co-speech movements all exhibit entrainment 
to underlying cortical cycles. Syntax itself can be interpreted as a rhythmic parsing problem: 
phrases unfold in ordered segments, constrained by oscillatory timing.  
Cross-linguistic studies reveal systematic variation in rhythm (stress-timed vs syllable-timed languages), 
suggesting that cultural differences in language mirror differences in gait-like temporal organization.  
From the RSVP--CPG perspective, linguistic structure is a specialization of rhythmic sequencing already present 
in cortical CPG dynamics.  

\subsection{Social Cognition and Collective Gait}
Interpersonal synchrony is a robust predictor of empathy and affiliation.  
Marching, dancing, and collective singing generate entrained rhythms across individuals, 
producing a sense of shared identity. Neuroscience confirms that inter-brain coupling emerges 
in such contexts, often measured as phase-locking between participants.  
RSVP--CPG interprets this as an extension of cortical gait into a social domain: 
when two or more agents synchronize their cycles, they create a temporary joint trajectory 
that scaffolds cooperation. Cultural variation in collective rhythms --- from ritual drumming 
to military drills --- can be seen as extensions of this basic neurodynamic principle.  

\subsection{Learning and Memory as Rhythm Stabilization}
Learning involves the creation of new gaits.  
Early skill acquisition is marked by irregular and unstable rhythms; with practice, 
these rhythms stabilize into reproducible cycles. Memory consolidation may be understood 
as the embedding of new rhythmic loops into the cortical repertoire.  
Disorders of memory, such as amnesia or dementia, can be reconceptualized as disruptions 
in the stability of these loops. Expertise, conversely, represents a refinement of gait: 
cycles become smoother, more efficient, and more entrainable with other processes.  

\subsection{Evolutionary Biology and Comparative Gait}
From an evolutionary standpoint, gait is one of the most conserved features across species.  
Neural central pattern generators controlling locomotion predate the neocortex by hundreds of millions of years.  
RSVP--CPG suggests that cortical cycles co-opted this ancient mechanism for higher cognition.  
Comparative research across species reveals both continuity and divergence:  
rodent hippocampal theta oscillations, avian song rhythms, and primate gestural repertoires 
all demonstrate gait-like dynamics.  
Genetic studies further indicate that core oscillatory parameters --- e.g.\ clock genes regulating circadian rhythms --- 
are deeply conserved. Thus RSVP--CPG may be framed as an evolutionary extension of motor control principles 
to the domain of conscious cognition.  

\subsection{Applied Connections}
The gait framework also informs applied sciences.  
In human-computer interaction, skeuomorphic design works precisely because it entrains novel tools 
to pre-existing motor gaits. In education, rhythmic scaffolding (e.g.\ musical training, patterned recall) 
enhances learning. In clinical therapy, interventions such as dance therapy, rhythmic auditory stimulation 
for Parkinson’s patients, and group synchrony exercises for autism spectrum disorders all leverage 
the brain’s preference for cyclical entrainment. RSVP--CPG provides a unifying rationale for these interventions.  

\subsection{Synthesis}
Across language, social life, learning, evolution, and application, cognition emerges as rhythm.  
The RSVP--CPG framework therefore not only contributes to theoretical neuroscience, 
but also offers a bridge across disciplines: from anthropology to education, 
from genetics to clinical medicine, all domains converge on the recognition 
that consciousness is not static representation, but ordered oscillation --- a gait.  

\newpage
\bibliographystyle{apalike}
\bibliography{references}

\end{document}
